\subsection{Degree-loss and remote-root compactifications}
\label{subsec:deg-degree-drop}

At degree loss the three equatorial roots occupy two scales, and the two
incident saddle exponents vanish at different rates. We construct one
through interval before dividing it according to accumulated action.
This ensures that the local estimates can later be combined for the
complete physical displacement.

\begin{proposition}[Degree-loss passage preparation]\label{deg:prop-degree-drop}
Let \(\delta=\kappa\tau\), with \(0<\kappa,\tau\ll1\), and assume
\begin{equation}
 \begin{gathered}
 a=\kappa A_{\rm d},\qquad
 \mathcal R a c=B\kappa,\qquad
 a+c=1-3\nu,\\
 \mathcal R(1-3\nu)=\omega+B\kappa .
 \end{gathered}
 \label{deg:degree-graph}
\end{equation}
The quantities
\[
 B,\ \omega,\ \mathcal R,\ A_{\rm d},\ c,\ 1-3\nu,\ c-a
\]
and their reciprocals are required to lie in prescribed compact
intervals.  Consider the source coefficients
\begin{equation}
 (a_s,p_s,\gamma_s,s_s)=
 \left(-1,\sqrt\tau\frac{\omega}{\sqrt B},
       \tau\kappa\mathcal R\nu,
       \tau^{3/2}\kappa^2\mathcal R\sqrt B\right).
 \label{deg:degree-source}
\end{equation}
On the fixed physical passage, below the common raw-height bound, the
ordered through interval is covered by the bounded-source, growing-action,
finite-middle-height, and fixed-exit ranges constructed below, together
with their high-height continuation. The actual matching maps,
derivative estimates, and strict overlaps are uniform on the displayed
compact coefficient family for positive \(\kappa,\tau\). At
\(\kappa\tau=0\), only the one-sided limits of normalized estimates
are asserted. The finite-\(V\) and fixed-exit ranges retain the original
section maps and finite closing on their actual buffered domains.
\end{proposition}

\begin{proof}
Set
\[
 r_s=\delta\mathcal R,\qquad
 v-\gamma_s=r_su,\qquad z=r_s^2w,\qquad
 \gamma_s=r_s\nu .
\]
The full two-dimensional system, after a positive change of time, is
\begin{equation}
 \dot u=\delta p_d(u)-K_d(u)w,\qquad
 \dot w=-w\{L_d(u)+E_d(u)w\},
 \label{deg:degree-field}
\end{equation}
where \(p_d=p_{\rm d}\) from \eqref{fx:pd} and
\((K_d,L_d,E_d)=(K,L,E)\) from \eqref{fx:coefficients}.
Thus the saddle rates and exponents are exactly those in
\eqref{fx:exponents}. On the present coefficient family,
\begin{equation}
 \chi=\frac{\delta\ell_+}{b_+}=\tau X_d,\qquad
 \beta=\frac{\delta\ell_-}{b_-}=\kappa\tau B_d,\qquad
 \frac{\beta}{\chi}=\kappa Q_d,
 \label{deg:degree-exponents}
\end{equation}
with \(X_d,B_d,Q_d\) positive units.  These identities provide the
uniform spectral gaps used in the two local saddle charts.

Fix \(0<\eta_\pm<1\), set
\(u_+=a(1-\eta_+)\), \(u_-=-1+\eta_-\), and on
\([u_-,u_+]\) write
\[
 h=-\delta p_d
   =\delta\mathcal R(u+1)(a-u)(c-u)>0,\qquad
 A_m={L_d\over h}.
\]
The scalar equation and its partial fractions are
\eqref{eq:comparison-middle} and \eqref{fx:partial-fractions}.
For \(J(u)=\exp\int_{u_+}^uA_m\) and \(J_-=J(u_-)\),
\begin{equation}
 J(u)=
 \left({1+u_+\over1+u}\right)^{1/\beta}
 \left({a-u_+\over a-u}\right)^{1/\chi}
 \left({c-u\over c-u_+}\right)^{1+1/\beta+1/\chi}.
 \label{deg:degree-middle-cocycle}
\end{equation}
Writing
\[
 \Lambda_\beta=
 \log{(1+u_+)(c-u_-)\over(1+u_-)(c-u_+)},
 \qquad \Lambda_\chi=\log a+O(1),
\]
we obtain
\begin{equation}
 \log J_-={\Lambda_\beta+(\beta/\chi)\Lambda_\chi\over\beta}
             +O(1),\qquad J_-\ge e^{c_0/\beta}.
 \label{deg:degree-large-cocycle}
\end{equation}
Indeed \(\Lambda_\beta\ge c_0>0\), while
\((\beta/\chi)|\log a|=O(\kappa|\log\kappa|)=o(1)\).

To measure how the middle passage concentrates near its negative end,
use the normalized action \(\rho=h(u_-)J/(J_-h)\) of
\eqref{fx:rho-gamma}, so
\(\rho(u_-)=1\).
If \(D_m=L_d-h'\), direct differentiation gives
\[
 {d\rho\over\rho}={D_m\over h}\,du,\qquad
 D_m'(u)=1-4\delta\mathcal R(u+\nu)>0.
\]
For one fixed \(u_0\in(u_-,0)\), this implies
\(\rho(u_0)\le e^{-c/\beta}\) and
\(\operatorname {Var}\rho\le2\).  Since
\[
 {h'\over h}={1\over u+1}-{1\over a-u}-{1\over c-u},
\]
the explicit expression \eqref{deg:degree-middle-cocycle} also yields
\begin{equation}
 \int_{u_-}^{u_+}\rho
 \left(|E_d|+\left|{h'\over h}\right|\right)du=o(1).
 \label{deg:degree-small-mass}
\end{equation}
The only noncompact term is bounded by
\(C\beta+C(1+|\log a|)e^{-c/\beta}/a\).  Thus no quotient
\(\nu/\kappa\) has entered the estimate.

We next place this middle transition inside a single interval of
trajectories joining the fixed physical sections.  Put
\begin{equation}
 s_0=\kappa\sqrt B,\qquad
 x_*={a\over s_0}={A_d\over\sqrt B},\qquad
 u=s_0\xi,\qquad w=s_0^2\tau y.
 \label{deg:degree-parabolic-dictionary}
\end{equation}
With
\[
 \Lambda_1=\frac{\omega}{\sqrt B},\qquad
 \Lambda_2=\kappa^2\mathcal R\sqrt B,\qquad
 \Lambda_3=\kappa\mathcal R\nu,
\]
the exact axis polynomial is
\begin{equation}
 P_s(\xi)=(1+s_0\xi)
       \left(1-{\xi\over x_*}\right)
       \left(1-{s_0\xi\over c}\right),
 \label{deg:degree-source-polynomial}
\end{equation}
and the full source system has the form
\begin{equation}
 \dot\xi=-\tau\{P_s+Q_\tau y\},\qquad
 \dot y=-y\{\xi B_\tau+C\tau^3\Lambda_2
                   (\Lambda_3+\Lambda_2\xi)y\},
 \label{deg:degree-source-field}
\end{equation}
where
\[
 B_\tau=1+\tau(2\Lambda_3+\Lambda_2\xi),\qquad
 Q_\tau=1-A\tau(\Lambda_3+\Lambda_2\xi)
          +C\tau^2(\Lambda_3+\Lambda_2\xi)^2.
\]

We construct two boundaries of the relevant through interval on the
same fixed physical sections.  For the upper comparison, prescribe the
actual point \(h_L=0\) of \eqref{bd:incoming-phase}, with
\(\eta=\tau y\).  Put
\(b_\tau=-\log(2\tau)\), \(q_\tau=b_\tau^{-1/2}\).
Freezing the source height scale changes the active quartet to
\begin{equation}
 \left(-q_\tau^2,\sqrt\tau\Lambda_1q_\tau,
       \tau\Lambda_3,{\tau^{3/2}\Lambda_2\over q_\tau}\right)
       \longrightarrow0.
 \label{deg:degree-zero-source}
\end{equation}
Lemma~\ref{bd:lem-source-ranges} constructs this orbit by an actual
incoming-to-central first hit, together with a real larger source
chord, using \eqref{bd:graph} with \(P_\sigma=P_s\).
Choose the fixed \(L>4(1+\sup x_*)\) of that lemma. Its energy
comparison gives, uniformly also on the larger source chord,
\begin{equation}
 \eta(\pm L)={L^2\over2}+O_L(\tau b_\tau),\qquad
 Q_\tau\eta+\tau P_s\ge c\eta .
 \label{deg:degree-zero-cuts}
\end{equation}
The same proof gives
\(\eta\asymp\tau b_\tau+\xi^2\) and the denominator margin
throughout the central source chord.  Thus \(h_L=0\) is an interior physical orbit, with neighbouring
through trajectories on both sides.

To continue to the fixed raw sections, put \(W=w/u^2\) and
\(r=|u|\).  Direct substitution in \eqref{deg:degree-field} gives
\begin{equation}
 rW_r={WA_\pm(r,W)\over
 K_d(\pm r)W-\delta p_d(\pm r)/r^2}-2W,
 \label{deg:degree-reciprocal}
\end{equation}
where
\[
 A_+={L_d(r)+E_d(r)r^2W\over r},\qquad
 A_-=-{L_d(-r)+E_d(-r)r^2W\over r}.
\]
After choosing a fixed \(U>2\max(1,c,|\nu|)+2\), then \(L\), the
strip
\begin{equation}
 {1\over3}<W<{3\over4},\qquad
 K_d(\pm r)W-\delta p_d(\pm r)/r^2>c_1
 \label{deg:degree-parabolic-strip}
\end{equation}
is invariant from \(s_0L\) to \(U\).  Equation
\eqref{deg:degree-zero-cuts} enters its interior.  The same reciprocal
equation in raw coordinates continues it to the fixed sections and
crosses each intervening equatorial root at positive height.

The lower boundary is the broken trajectory
\begin{equation}
 W^s_+(a,0)\cup\{(u,0):-1\le u\le a\}\cup W^u_+(-1,0).
 \label{deg:degree-stopped-curve}
\end{equation}
The invariant-manifold theorem in the two fixed weighted half-boxes
constructs the upper stable and unstable branches.  Their graph
denominators have a positive margin after division by the local
transverse scale.  To attach them to the fixed sections, first leave
the positive box at \(\xi=x_*+\ell_*\), with \(\ell_*>0\) fixed.
Its height and denominator are bounded below by \(c\ell_*\).
The graph equation therefore reaches \(\xi=L\), with
\[
 \eta_+^{\rm st}(L)=(L^2-x_*^2)/2+o(1).
\]
On the negative side leave the weighted box at \(u=-R\), with
\(R=1+\rho_*\), \(\rho_*>0\) fixed.  Its outgoing graph has
\[
 w_-^{\rm st}(-R)=(R^2-1)/2+o(1),\qquad
 K_dw_-^{\rm st}-\delta p_d\ge c(R^2-1).
\]
The barriers \(2(r^2-1)/5\) and \(3(r^2-1)/5\) continue this
graph until it enters \eqref{deg:degree-parabolic-strip}.  Both
branches then reach \(u=\pm U\), and afterwards the fixed raw
sections.  The source cuts \(u=\pm s_0L\) occur before all three equatorial
roots have been passed. The intermediate fixed cuts complete that
continuation before the raw tails are attached.

Between the saddle boxes \(p_d<0\), so
\(\dot u=-(h+K_dw)<0\) for every \(w\ge0\).  On either exterior
side the stopped branch has \(\dot u<0\) away from its saddle and
\(\partial_w\dot u=-K_d<0\).  Scalar order thus keeps every graph
between the stopped and upper comparison graphs in the same positive
denominator domain.  The inclination lemma in the two saddle boxes,
together with Lemma~\ref{lem:loc-through}, identifies the stable
branch as the lower endpoint of the through interval.  Write the two
actual positive-section traces as
\(h_{\rm in}^{\rm st}<h_{\rm in}^0\), where \(h_{\rm in}^0\) is the
trace of \(h_L=0\). There is no cycle at or below
\(h_{\rm in}^{\rm st}\).

The quadratic outer equation contracts towards its organizing corridor
in logarithmic radius.  More explicitly, for a fixed small raw cut
\(v_0\), its deviation at \(\pm v_0\) is bounded by
\[
 C(r_sU/v_0)^\alpha+Cr_s/v_0
       +C(|A-1|+|C-1|),\qquad 1<\alpha<2.
\]
This follows from \(\mu_i=O(r_s^{4-i})\), \(i=1,2,3\), and
variation of constants after \(u=\pm U\), retaining the passive
coefficient differences from the organizing equation. The fixed regular maps
then put both boundary traces, and by scalar order the whole interval
\((h_{\rm in}^{\rm st},h_{\rm in}^0)\), in the prescribed raw section
windows.
This scalar order locates the orbits in the raw windows. We estimate
the derivatives separately on the enlarged chords transported into
those windows.

The two fixed comparisons of Proposition~\ref{bd:prop-bounded-action}
supply the overlapping source interval $-\delta_2\le h_L\le0$ and
bounded strong-footpoint annulus. Their enlarged independent windows
reach the same raw sections by \eqref{bd:parabolic-cuts} and the outer
bridges above.

Let \(r_b\) be the negative strong footpoint where that section is
met.  On the lower part of the through interval let \(r_\ell\) be
the first point on the actual curve \(y/P_s=\ell_0\).
Choose a fixed \(r_c>2\).  A fixed real comparison through orbit
has a positive ratio at \(-r_c\); choose \(\ell_0\) below that
ratio.  Scalar order gives a unique lower interval on which the
initial ratio is smaller than \(\ell_0\).  Before its first hit,
directly from \eqref{deg:degree-source-field},
\begin{equation}
 {d\over dr}\log{y(-r)\over P_s(-r)}
 \ge {c\over\tau(1-s_0r)}
 \label{deg:degree-ratio-monotonicity}
\end{equation}
on the full interval \(r_c\le r<1/s_0\).  A strict through graph
has positive height at the negative root while \(P_s(-r)\to0\).
Hence the hit exists.  Its uniqueness and the positive first
variation at a fixed section give
\begin{equation}\label{deg:degree-phase-order}
 \partial_{h_{\rm in}}r_\ell<0,\qquad
 \partial_{h_{\rm in}}r_b<0,\qquad
 \frac{dr_\ell}{dr_b}>0
\end{equation}
on the common domain, with \(h_{\rm in}\) the actual raw input height.
No equality between the two footpoint coordinates has been used.

We need a coefficient-independent threshold at which a first ratio
hit has sufficiently large action.  Fix a positive weak-section input \(P_0\) in
the positive saddle box.  Its local Dulac formula and
\(\chi\asymp\tau\) give
\(-\tau\log z_+\le C(P_0)\) for inputs \(P\ge P_0\).
The exact core equation consequently gives
\(\log(y/P_s)(-r_c)\ge-C(P_0)/\tau\), unless the first hit
has already occurred.  If there were no hit before \(R\),
\eqref{deg:degree-ratio-monotonicity} would imply
\[
 \log(y/P_s)(-R)\ge
             \{-C(P_0)+c(R-r_c)\}/\tau.
\]
Choosing \(R=R(P_0)\) large is a contradiction.  Thus
\(r_\ell\to\infty\) forces \(P\to0\) uniformly.  Choose
\(P_0\) so that its positive action exceeds the fixed middle
variational exponent by two; then fix
\(r_{\rm ge}^{\rm out}>R(P_0)\) and
\(r_{\rm ge}>r_{\rm ge}^{\rm out}\).  The growing-action estimates below therefore start with a strict lower
action margin.

At \(r_\ell=r_{\rm ge}\), the actual strong projection has a
compact image.  Choose an independent \(R_{\rm BD}\) so that
this image lies strictly below \(2R_{\rm BD}\).  The transverse
inverse and \eqref{deg:degree-phase-order} then give a fixed
\(r_{\rm ge}^{\sharp}>r_{\rm ge}\) whose entire closed interval
\([r_{\rm ge},r_{\rm ge}^{\sharp}]\) lies inside the enlarged
annular resection.  In particular,
\[
 r_b>2R_{\rm BD}\ \Longrightarrow\ r_\ell>r_{\rm ge}.
\]
Only now decrease the upper \(s_0\)-bound so that this whole fixed
overlap lies in the growing-action product range.  Fix \(V_1>1\) and
\(0<\sigma_1<\sigma_2<\sigma_0<1/4\), with
strictly enlarged ranges at every endpoint.
The four passage regions will be restrictions of this ordered through
interval.

Choose one coefficient-only real action offset \(A_{c,0}(\theta)\),
held fixed by every phase derivative, and put
\[
 \mathcal P(r)=P_s(-r)/B_\tau(-r),\qquad
 \mathcal D=(\mathcal P(r)/r)\partial_r,\qquad
 \mathcal A_-(r)=\int_{r_c}^r\frac{vB_\tau(-v)}{P_s(-v)}\,dv-A_{c,0}.
\]
Lemma~\ref{deg:lem-growing-matching-direct} constructs the two
centre maps independently, retaining the exact initial polynomial
and both core multipliers. Its product rows
\eqref{deg:source-growing-product-phase-rows} and matched identity
\eqref{deg:degree-growing-matching} give the actual graph, with
\(|\zeta|+|X\partial_X\zeta|\le C\tau\) on its enlarged image.

Construct one negative section map for both the growing and fixed-exit
ranges. Fix \(u_c\in(-1,u_-)\) with \(|u_c|>2\sigma_0\), retain
its enlarged product buffer \(\sigma_0\le\bar\sigma<|u_c|/2\), and choose
\(\Delta>0\) below one third of the minimum of the gap
\eqref{fx:overlap-gap} on \([\sigma_1,\sigma_2]\), also with
\(\Delta<\frac14\inf_c\Gamma_{-,c}\). Enlarge the intervals of
\eqref{fx:W-limit} once more, with uniform strict margins:
\[
 \mathcal J_c\Subset\mathcal J_c^\sharp
 \Subset\mathcal J_c^{\sharp\sharp}\Subset(0,u_c^2/2).
\]
Apply Lemma~\ref{deg:lem-growing-landing-direct} on the original
enlarged lip range, with its fixed lower threshold and product buffer.
Equations \eqref{deg:landing-direct-value} and
\eqref{deg:landing-direct-gamma-unit} give the actual landing value,
its strict derivative and its positive first-variation unit.
The action supremum in \eqref{deg:landing-direct-first-row} is
$O(\tau(1+|\log\tau|))$, uniformly in independent $\kappa,j$.
Since \(s_0r\le\sigma_0\), these heights lie strictly inside
\[
 \mathcal J_{\rm gr}=[u_c^2/4,3u_c^2/4]
 \Subset\mathcal J_{\rm gr}^{\sharp}=(u_c^2/8,7u_c^2/8).
\]
Choose a fixed compact positive-height interval \(\mathcal I_{\rm out}\)
containing the closures of \(\mathcal J_{\rm gr}^{\sharp}\) and
\(\mathcal J_c^{\sharp\sharp}\) with a strict buffer, uniformly in
\(c\). Thus it contains \(u_c^2/2\), while the fixed-exit chord
stays strictly below that height.

For every \(W\in\mathcal I_{\rm out}\), the limiting graph
\(w(u)=W+(u^2-u_c^2)/2\) has a uniform positive height from
\(u_c\) to a fixed large negative cut \(-U_-\). Ordinary ODE
dependence gives a positive exact denominator, including at \(u=-1\),
and uniform \(C^5\) bounds for this finite passage and its inverse.
Choose \(U_-\) to enter \eqref{deg:degree-parabolic-strip}.
The tail estimates
\eqref{ann:tail-first-variation}--\eqref{eq:N-parameter-bound},
with the exact first cocycle extracted through three variations, give
\[
 \begin{gathered}
 z_-(W)=z_-^{\rm st}+\Omega Z_-(W),\qquad
 \Omega=r_s^{\,2-\lambda_\gamma},\quad
 \lambda_\gamma=\frac{12\gamma_s}{1+2\gamma_s},\\
 Z_-'>c,\qquad \|Z_-\|_{C^3}+\|Z_-^{-1}\|_{C^3}\le C.
 \end{gathered}
\]
Since \(\Omega\to0\) and the stopped trace lies inside its fixed raw
window, the whole chord is in the domain of \(d_\theta\) after shrinking
the coefficients. Compose with that map, put \(t_b=d_\theta(z_-^{\rm st})\),
and fix an interior reference \(W_\circ\). With
\(U_{\rm out}=\Omega^{-1}(d_\theta\circ z_-)'(W_\circ)\),
this defines a single actual map
\begin{equation}\label{deg:fixed-exit-actual-output}
 \begin{gathered}
 \Xi(W)=t_b+\Omega U_{\rm out}D^n(W),\qquad U_{\rm out}^{\pm1}\le C,\\
 (D^n)'>c,\qquad \|D^n\|_{C^3}+\|(D^n)^{-1}\|_{C^3}\le C.
 \end{gathered}
\end{equation}
For compact \(\nu\), \(|\lambda_\gamma\log r_s|\to0\), so
\(\Omega\asymp r_s^2\); the exact power is retained for the remote-root
face. In particular \(\Xi'\asymp\Omega\) and
\(|\Xi''/\Xi'|\le C\) on the entire common chord.

Fix \(r_\circ\) inside the enlarged growing domain and define
\[
 \begin{gathered}
 a_{\rm grow}=\Omega s_0^2U_{\rm out},\qquad
 t_\dagger=\Xi(W_\theta(r_\circ))+a_{\rm grow}r_\circ^2/2,\\
 T(r)=\Xi(W_\theta(r))=t_\dagger-a_{\rm grow}\mathcal B(r).
 \end{gathered}
\]
The fixed-endpoint variation in
Lemma~\ref{deg:lem-growing-landing-direct} gives
\begin{equation}\label{deg:degree-growing-output}
 \mathcal B'(r)=rG(r),\qquad
 G(r)=\Gamma_\theta(r)
              \frac{\Xi'(W_\theta(r))}{\Omega U_{\rm out}}>0.
\end{equation}
Its exact chain rule is
\begin{equation}\label{deg:source-growing-output-log-row}
 \mathcal D\log G=
 \mathcal D\log\Gamma_\theta
 -s_0^2\mathcal P(r)\Gamma_\theta(r)
                  \frac{\Xi''(W_\theta(r))}{\Xi'(W_\theta(r))}.
\end{equation}
Estimates~\eqref{deg:landing-direct-raw-first-row} and
\eqref{deg:source-growing-actual-quotient-estimates} control the two
terms that enter the complete displacement after finite closing.

On the positive side use the fixed local cuts and stopped outer chord
of Lemma~\ref{lem:actual-dulac-matching}. Its backward-integral and tail
construction holds through five phase derivatives. Thus the natural
raw input is
\begin{equation}\label{deg:fixed-exit-natural-input}
 P=\frac{x_{\rm raw}-x_b}{s_{\rm in}},\qquad
 s_{\rm in}=\Omega a^{2-\lambda_\gamma}U_{\rm in},\qquad
 U_{\rm in}^{\pm1}\le C,
\end{equation}
and the normalized map from the fixed outer chord and its inverse
have uniform \(C^5\) bounds. In the positive saddle coordinates
\(x=(u-a)/a,y=w/a^2\), write the invariant restrictions on
\(x=-\eta_+\) as
\[
 A_+(y)=-\Xi_{+,m}(y),\quad A_0=A_+(0)>0,\qquad
 Z_{+,m}(y)=z_+y g_+(y),\quad g_+(0)=1.
\]
Here \(A_0,z_+\) are positive units and \(A_+,g_+\) have uniform
finite derivative bounds. On the outer chord the invariant coordinate
constructed in that lemma is
\(\mathscr L_\theta(P)=-I(P)/A_0=\Phi_+(P)/A_0\); hence
\[
 \mathscr L_\theta(0)=0,\qquad \mathscr L_\theta'>c>0,\qquad
 \|\mathscr L_\theta\|_{C^5}
       +\|\mathscr L_\theta^{-1}\|_{C^5}\le C.
\]
Choose \(p_0>0\) small, independently of the coefficients, so that
\(\mathscr L_\theta(P)\le\vartheta<1\) on \([0,p_0]\).
Equality of the two invariants is exactly
\begin{equation}\label{deg:fixed-exit-middle-invariant}
 \mathscr L_\theta(P)=\frac{A_+(y)}{A_0}
                         \{z_+y g_+(y)\}^{\chi}.
\end{equation}
This equation has a unique small nonnegative solution \(y=y(P)\).
Indeed, on \(0<y\le\chi^2\) the logarithmic derivative of its
right side, multiplied by \(y\), is
\[
 B_\chi(y)=\chi+y\frac{A_+'(y)}{A_+(y)}
                   +\chi y\frac{g_+'(y)}{g_+(y)}
          =\chi+O(y)\ge\chi/2.
\]
The right side of \eqref{deg:fixed-exit-middle-invariant} tends
uniformly to one at \(y=\chi^2\), while its limit at zero is zero.
Thus it covers \([0,\vartheta]\) for
all sufficiently small \(\chi\). Taking logarithms gives
\[
 y=z_+^{-1}\mathscr L_\theta(P)^{1/\chi}
       \exp O(y/\chi+y).
\]
In particular the exponential factor is a positive unit. From
\(dy/d\mathscr L=y/(\mathscr L B_\chi(y))\), differentiation
twice more and the bounded derivatives of \(A_+,g_+\) give
\begin{equation}\label{deg:fixed-exit-flat-height}
 |\partial_P^j y(P)|\le
 C_j\chi^{-j}\mathscr L_\theta(P)^{1/\chi-j},
 \qquad 0\le j\le3,\quad P>0.
\end{equation}
For \(\chi<1/4\) these derivatives have zero limits at \(P=0\).
Their suprema on the fixed interval tend to zero, since
\(\chi^{-j}\vartheta^{1/\chi-j}\to0\).

Define, with its continuous value at zero,
\begin{equation}\label{deg:fixed-exit-exact-weak-coordinate}
 \Phi_\theta(P)=\{z_+y(P)\}^{\chi}
 =\mathscr L_\theta(P)\frac{A_0}{A_+(y(P))}
                          g_+(y(P))^{-\chi}.
\end{equation}
The second expression and \eqref{deg:fixed-exit-flat-height} show
that \(\Phi_\theta-\mathscr L_\theta\to0\) in \(C^3\).
Consequently \(\Phi_\theta'>c>0\), its \(C^3\) norm and inverse
norm are uniform, and \(\Phi_\theta(P)=PU_\theta(P)\) with
\(U_\theta\) a positive unit on \([0,p_0]\).
Only phase derivatives are used here.

Use \(\lambda=\ell_0\) in the normalized middle action
\(w=\ell_0h\rho V\) of \eqref{eq:comparison-action}.
The exact reference value follows directly from that raw height. Put
\[
 \varepsilon_+=\frac{\lambda h_-}{a^2J_-},\qquad
 \Lambda_+^m=(z_+\beta\varepsilon_+)^\chi,
 \qquad
 \Phi_\theta(a_{\rm in})=\Lambda_+^m.
\]
Write \(q_s=\beta/\chi\), \(j=1/c\), \(d=\mathcal Rc\), and
\(\mathscr R=(1-ju_-)/(1-ju_+)\). The logarithms in
\eqref{deg:degree-middle-cocycle} are exactly
\[
 \Lambda_\beta=\log\frac{(1+u_+)\mathscr R}{1+u_-},\qquad
 \Lambda_\chi=\log\frac{(a-u_+)\mathscr R}{a-u_-}.
\]
The reference telescope is
\begin{equation}\label{p2:reference-telescope}
 \begin{gathered}
 \beta\log(z_+\beta\varepsilon_+)
 =-\Lambda_\beta-q_s\Lambda_\chi+\beta\log(\tau^2U_{\rm tel}),\\
 U_{\rm tel}=\frac{z_+\ell_0\beta\delta d}{a^2\tau^2}
                (1+u_-)(a-u_-)(1-ju_+),\qquad U_{\rm tel}^{\pm1}\le C.
 \end{gathered}
\end{equation}
Here \(\delta=\kappa\tau\), \(a=\kappa A_d\), and
\(\beta=\kappa\tau B_d\); hence all factors of \(\kappa\)
cancel in \(U_{\rm tel}\). Since
\(\Lambda_\beta=\Gamma_{-,c}+O(a)\) and
\(\Lambda_\chi=\log a+O(1)\), the error is bounded by
\(C\{a+q_s(1+|\log a|)+\beta(1+|\log\tau|)\}\).
It tends to zero under independent coefficient bounds, giving
\begin{equation}\label{deg:degree-fixed-exit-gap}
 q_s\log\Phi_\theta(a_{\rm in})
 =\beta\log(z_+\beta\varepsilon_+)
 =-\Gamma_{-,c}+o(1).
\end{equation}
Thus \(\Lambda_+^m\to0\) and the reference
\(a_{\rm in}\in(0,p_0)\) exists uniquely. It corresponds to
\(y=\beta\varepsilon_+\), hence to the raw middle input \(V=\beta\).
For every orbit in this local input interval, the exact height dictionary is
\begin{equation}\label{fx:height-dictionary}
 \begin{gathered}
 y=\varepsilon_+V,\qquad
 V=\beta\left(\frac{\Phi_\theta(P)}
                    {\Phi_\theta(a_{\rm in})}\right)^{1/\chi},\\
 \frac{w(u_+)}{h(u_+)}=\lambda\rho(u_+)V.
 \end{gathered}
\end{equation}

We have constructed the through interval, its first ratio hits and its
input, output and centre-height maps. Proposition~\ref{p2:prop-unified-carrier-cover}
uses them to cover every input, including truncated intervals and
boundaries between the four ranges. Proposition~\ref{p2:prop-finite-v-fe-actual-budget}
constructs the last two transitions on the same shooting chord and
buffered closing image.
\end{proof}

The remote positive root requires a different coordinate at infinity.
Its reciprocal keeps the finite passage coefficients regular while the
physical tail scales are retained exactly.

\begin{proposition}[The remote-root face]\label{deg:prop-degree-infinity}
Let \(d,A_d\) range in positive compact intervals and put
\begin{equation}
 \begin{gathered}
 0\le j\le j_0,\qquad a=\kappa A_d,\qquad
 \mathcal R=dj,\qquad B=dA_d,\\
 \omega=d(1-a+ja),\qquad
 j\nu={j-ja-1\over3},\qquad \delta=\kappa\tau .
 \end{gathered}
 \label{deg:degree-infinity-graph}
\end{equation}
For \(j>0\) this is \eqref{deg:degree-graph} with \(c=j^{-1}\).
For \(j>0\), the conclusion of Proposition~\ref{deg:prop-degree-drop}
extends uniformly as \(j\downarrow0\). At \(j=0\) this denotes only the
one-sided closure of the analytic estimates, not an additional physical
coefficient fibre or an asserted inverse raw gate.
\end{proposition}

\begin{proof}
The factorization and translated root are
\begin{equation}
 p_d(u)=d(u+1)(u-a)(1-ju),\qquad
 \gamma_s={\kappa\tau d\over3}(j-ja-1).
 \label{deg:degree-infinity-factor}
\end{equation}
All coefficients of the exact field on bounded \(u\)-intervals are
analytic at \(j=0\).  In the middle calculation the apparently singular
quantities reduce to
\[
 {1-ju\over1+j},\qquad {j\over1-ju},\qquad
 D_m'=1-4\gamma_s-4\kappa\tau dju.
\]
The limiting functions needed in the accumulated-action and fixed-exit
arguments are
\[
 H_{0,j}(u)=-{u(1+u)(1-ju)\over1+j},\qquad
 \Gamma_j(u)=
 \log{(1+u)(1-ju_-)\over(1+u_-)(1-ju)},
\]
and
\[
 \Gamma_{-,j}-\Gamma_j(-\sigma)
 =\log{1+j\sigma\over1-\sigma}>0.
\]
Thus the small-mass, Lambert, and fixed-exit estimates above extend
continuously to \(j=0\).

The bounded-action domains also have regular coordinates.  Put
\(\sigma=s_0=\kappa\sqrt B\), \(x_*=A_d/\sqrt B\).  Then
\[
 P_\sigma(\xi)=(1+\sigma\xi)(1-\xi/x_*)(1-\sigma j\xi),\qquad
 \Lambda_2=\sigma^2j/x_*,\quad
 3\Lambda_3=\sigma(j-1)/x_*-\sigma^2j<0.
\]
The positive root remains simple and the negative compact-output
interval has a positive polynomial margin.  The two fixed phase
comparisons, the compact central-height estimate, and the negative
frozen first-exit construction in
Lemma~\ref{bd:lem-source-ranges} are consequently uniform on the
compact \(j\)-block.  Lemma~\ref{bd:lem-annulus} uses the sign
\(\Lambda_3<0\) on this face: its actual first tail cocycle is only
bounded above by \(Ce^2\), with \(e=\tau\Lambda_2\), and is
retained exactly in the relative normalization.  Thus it requires
no bound on \(\tau|\Lambda_3||\log e|\), even when \(j\) is
arbitrarily small.  The bounded-action estimate is therefore uniform even as the third root
escapes the compact root chart.

It remains to check the unbounded tails.  Their exact source coefficients
are
\[
\begin{array}{c|cccc}
 &a_\pm&p_\pm&\gamma_\pm&s_\pm\\ \hline
 +&-\tau d/A_d&\tau d(1-a+ja)/A_d&\gamma_s&r_sa\\
 -&-\kappa^2\tau dA_d&\kappa\tau d(1-a+ja)&\gamma_s&r_s,
\end{array}
\qquad r_s=\kappa\tau dj.
\]
After the transverse blow-down, the two maps to the fixed raw sections
have first-cocycle scales
\begin{equation}
 r_s^{\,2-\lambda_\gamma}U_-,\qquad
 r_s^{\,2-\lambda_\gamma}a^{\,2-\lambda_\gamma}U_+,
 \qquad \lambda_\gamma={12\gamma_s\over1+2\gamma_s},
 \label{deg:degree-common-scale}
\end{equation}
with normalized phase derivatives bounded.  The common factor
\(r_s^{2-\lambda_\gamma}\) cancels in the complete transition, while
\(|\lambda_\gamma\log a|=o(1)\). With \(q_s=\beta/\chi\),
the exact telescope \eqref{p2:reference-telescope} gives
\[
 q_s\log\Lambda_+^{m}=-\Lambda_\beta+o(1),\qquad
 e^{-C/q_s}\le\Lambda_+^{m}\le e^{-c/q_s},
\]
and contains no term \(\beta\log j\). The common section scales are
\(s_{\rm in}=\Omega a^{2-\lambda_\gamma}U_{\rm in}\) and
\(a_{\rm out}=\Omega U_{\rm out}\), with
\(\Omega=r_s^{2-\lambda_\gamma}\).

The local boxes and normalized outer maps have uniform finite phase
bounds on the compact $j$-block, so the weak-coordinate construction
\eqref{deg:fixed-exit-middle-invariant}--\eqref{deg:fixed-exit-flat-height}
has exact reference
\[
 \Phi_\theta(a_{\rm in})=(z_+\beta\varepsilon_+)^\chi,
 \qquad q_s\log\Phi_\theta(a_{\rm in})=-\Gamma_{-,j}+o(1).
\]
Proposition~\ref{p2:prop-finite-v-fe-actual-budget} uses these regular
functions \(H_{0,j},\Gamma_j\), their strict action gap and the exact
raw scales to extend the last two passage estimates to this face.

For the outer geometry no fixed cut beyond \(|\nu|\) is available.
Choose instead the positive source cut \(\xi=x_*L_t\), so that
\(u/a=L_t\) exactly.  The actual source construction gives
\(1/3<\eta/\xi^2<3/4\) and a positive divided denominator there.
On the negative side continue from \(u=-s_0L\) only to
\(u=-L_t\).  On this interval the regular factorization gives
\[
 |p_d(-r)|\le Cr,\qquad
 |\delta p_d(-r)|/r^2\le C\tau/L,\qquad
 |K_d-1|+|-L_d(-r)/r-1|\le C\kappa\tau.
\]
The equation for \(w/u^2\) therefore has the same invariant strip
and positive denominator.  These are the two finite entry chords
for the bounded source quartets in the table.  Their normalized
tail equations \eqref{eq:parabolic-tail} continue to the fixed raw
sections; the positive tail crosses \(u=c=j^{-1}\), if that root
lies before the section, at positive height.  The stopped branches
are attached by the same weighted local graphs and finite bridges.
Since \(p_d<0\) on \(-1<u<a\), and \(\partial_w\dot u=-K_d<0\),
every intermediate graph remains on the through side of these
branches.  This constructs the entire ordered interval before its
zero bounds are applied.

The same ratio and strong-projection orders retain the bounded/growing
overlap. The positive $\Gamma_j$ gap retains the fixed-exit overlap.
On the real enlarged source chord through the actual cut \(h_L=0\),
\eqref{bd:positive-height} again gives
\(Y_0\asymp b_\tau\to\infty\), while
\(s_s^2Y_0\le C\tau^3\kappa^4j^2b_\tau\to0\).
Hence this chord supplies the overlap of the actual high interval
with the low carrier through $h_L=0$.
All four domains and estimates are uniform for $j>0$;
$j=0$ retains only the normalized one-sided meaning.
\end{proof}

\subsection{Direct negative landing and its first action variation}

A negative ratio hit must be continued to a fixed section before it can
be used in the physical return. The next lemma constructs this landing
and estimates its first action derivative. The terminal section is kept
fixed while the starting lip moves.

\begin{lemma}\label{deg:lem-growing-landing-direct}
Use the exact field \eqref{deg:degree-field} on either the compact
coefficient graph \eqref{deg:degree-graph} or the reciprocal graph
\eqref{deg:degree-infinity-graph}, with physical \(j>0\).
Keep the original sections, compact passive coefficients, and
coefficient-only scales. Put
\[
 \epsilon_d=\kappa\tau,\qquad s_0=\kappa\sqrt B,\qquad
 x_c=-u_c,\qquad \sigma=s_0r .
\]
Choose the original enlarged product buffer
\(0<\sigma_0\le\bar\sigma<x_c/2<1/2\), and fix
\(R_0>r_c\). On the original enlarged range
\begin{equation}\label{deg:landing-direct-domain}
 R\le r\le\bar\sigma/s_0,\qquad R\ge R_0,\qquad
 q_\ell=\frac{\epsilon_d}{\sigma}=\frac{\tau}{\sqrt B\,r},
 \qquad L_\ell=1+|\log q_\ell|,
\end{equation}
the graph from the original lip \(w=\ell_0h\) reaches \(u_c\)
in positive time, and its inverse reaches the original centre
\(u=0\). Both graphs have positive height and denominator.

Write \(W_\theta(r)=w(u_c;r)\). There is an exact positive unit
\(\Gamma_\theta(r)\) such that
\begin{equation}\label{deg:landing-direct-value}
 \left|W_\theta(r)-\frac{x_c^2-\sigma^2}{2}\right|
 \le C\epsilon_d\{1+\sigma L_\ell\},\qquad
 W_\theta'(r)=-s_0^2r\Gamma_\theta(r),
\end{equation}
and
\begin{equation}\label{deg:landing-direct-gamma-unit}
 c\le\Gamma_\theta\le C,\qquad
 |\Gamma_\theta-1|\le C(q_\ell L_\ell+\epsilon_d).
\end{equation}
If \(\Gamma(\sigma)=\Gamma_\theta(\sigma/s_0)\), then, with all
coefficients frozen,
\begin{equation}\label{deg:landing-direct-first-row}
 \begin{aligned}
 |\partial_\sigma\log\Gamma|
 &\le C\left\{\frac{q_\ell}{\sigma}L_\ell+\epsilon_d\right\},\\
 |\mathcal D\log\Gamma_\theta|
 &\le C\left\{\frac{q_\ell}{r}L_\ell+s_0\epsilon_d\right\},
 \qquad \mathcal D=\frac{\mathcal P(r)}r\partial_r .
 \end{aligned}
\end{equation}
Constants are fixed by the original compact ranges, \(R_0,r_c,\ell_0\)
and the section buffers, independently of the later \(R,\kappa,\tau,j\).
In particular the supremum of the action row is
\(O(\tau(1+|\log\tau|))\) as \(\tau\to0\), uniformly in independent
\(\kappa,j\). At fixed \(r,\tau>0\), sending only \(\kappa\to0\)
does not assert that this row tends to zero or that \(\Gamma_\theta\to1\).
\end{lemma}

\begin{proof}
Put \(x=-u\), \(H(x)=h(-x)\), \(k(x)=K(-x)\),
\(b(x)=-L(-x)\), and \(e(x)=-E(-x)\). The exact graph is
\begin{equation}\label{deg:landing-direct-graph}
 w_x=F(x,w)=\frac{w(b+ew)}{H+kw},\qquad
 w(\sigma;\sigma)=\ell_0H(\sigma).
\end{equation}
Indeed the original dictionary gives
\(s_0^2\tau P_s(-r)=h(-s_0r)\) exactly.
In the compact and reciprocal graphs, respectively,
\[
 H=\epsilon_d\mathcal R(1-x)(a+x)(c+x),\qquad
 H=\epsilon_dd(1-x)(a+x)(1+jx).
\]
On \([\sigma,x_c]\), \(a/\sigma\le C/R_0\).
The exact coefficients \eqref{fx:coefficients} therefore give
\begin{equation}\label{deg:landing-direct-coefficients}
 \begin{gathered}
 H\asymp\epsilon_dx,\quad |H'|+|H''|\le C\epsilon_d,\quad
 |H'/H|\le C/x,\qquad k\asymp1,\\
 |k-1|+|k'|\le C\epsilon_d,\qquad |b-x|\le C\epsilon_dx,\\
 |(b/x)'|\le C\epsilon_d,\quad
 |b'/b|\le C/x,\qquad |e|+|e'|\le C\epsilon_d^2 .
 \end{gathered}
\end{equation}
At the reciprocal face these estimates use
\(\epsilon_d\mathcal R=\epsilon_ddj\) and
\(\epsilon_d\mathcal R\nu=\gamma_s\); no bound on \(\nu\) or
\(\log j\) is required. A common small upper bound on \(\tau\)
makes \(q_\ell\) small without comparing any radial parameters.

First admit the graph. On a provisional band \(0<w<Mx^2\),
the denominator and \(b+ew\) are positive, and
\[
 0<F\le Cx+C\epsilon_d^2w,\qquad
 F\ge\frac{cxw}{C\epsilon_dx+Cw}.
\]
A fixed large \(M\) prevents first contact with \(Mx^2\);
the lip is below it since \(H(\sigma)/\sigma^2=O(q_\ell)\).
Thus the graph stays positive and reaches \(x_c\).
On \([\sigma,2\sigma]\), height is increasing and
\[
 \frac{dx}{dw}\le C\frac{\epsilon_d\sigma+w}{\sigma w}.
\]
For fixed small \(m>0\), the time to height \(m\sigma^2\) is at
most \(C\epsilon_dL_\ell+Cm\sigma<\sigma\), after a common
coefficient shrink. A smaller quadratic lower barrier then gives
\begin{equation}\label{deg:landing-direct-corridor}
 0<w\le Cx^2,\quad b+ew\asymp x,\qquad
 w\asymp x^2\quad(2\sigma\le x\le x_c).
\end{equation}

For the inverse to the centre, the same exact coefficients give,
on \([s_0r_c,\sigma]\) while \(w/H\le\ell_0\),
\[
 \frac{d}{dx}\log(w/H)
 =\frac{b+ew}{H+kw}-\frac{H'}H
 \ge c/\epsilon_d-C/x>0.
\]
Here \(\epsilon_d/x\le C\tau/r_c\). Backward continuation remains
below this lip, whose hit at \(\sigma\) is therefore the first hit
after the inner cut. At that cut its height \(w_b\) satisfies
\(w_b\le C\epsilon_d(a+s_0r_c)\le C\kappa^2\tau\).
On \([0,s_0r_c]\), \(H\ge c\epsilon_da>0\), and for
\(0<w\le2w_b\) the exact derivative satisfies
\[
 F\ge-C\frac{\epsilon_d^2}{\epsilon_da}w^2\ge-C\tau w^2.
\]
In the inverse direction this prevents first contact with \(2w_b\).
Positive denominator and scalar uniqueness give a positive centre
height. This argument permits either sign of \(e(0)w\).

Write \(D=H+kw\) and
\[
 N=\frac{H'+(k'+e)w}{D},\qquad
 I_x=\int_\sigma^xN(t,w(t;\sigma))\,dt,\qquad I=I_{x_c}.
\]
Direct differentiation along the graph gives
\(F_w=(\log(w/D))_x+N\).
The near height change and the quadratic corridor imply
\[
 \int_\sigma^{2\sigma}\frac{\epsilon_d}{D}\,dx
 \le Cq_\ell L_\ell,\qquad
 \int_{2\sigma}^{x_c}\frac{\epsilon_d}{D}\,dx\le Cq_\ell.
\]
The remaining part of \(N\) has integral \(O(\epsilon_d)\), so
\(|I_x|\le C(q_\ell L_\ell+\epsilon_d)\).
Keep the complete moving initial factor
\begin{equation}\label{deg:landing-direct-initial-factor}
 C_\sigma=(1+\ell_0k(\sigma))H'(\sigma)
                 -b(\sigma)-\ell_0e(\sigma)H(\sigma).
\end{equation}
The moving initial point contributes the intrinsic derivative
\(\ell_0H'(\sigma)-F(\sigma,\ell_0H(\sigma))\). Combining this
with the Jacobian to the fixed terminal section gives exactly
\begin{equation}\label{deg:landing-direct-exact-variation}
 w_\sigma(x;\sigma)=\frac{w(x;\sigma)}{D(x;\sigma)}C_\sigma e^{I_x},
 \qquad
 \Gamma(\sigma)=\frac{w(x_c;\sigma)}{D(x_c;\sigma)}
                       \frac{-C_\sigma}{\sigma}e^I .
\end{equation}
Here \(-C_\sigma/\sigma=1+O(q_\ell+\epsilon_d)\), and the terminal
ratio is \(1+O(\epsilon_d)\). These identities prove the strict sign,
the unit estimate, and \(|w_\sigma(x;\sigma)|\le C\sigma\).
The value estimate follows by integrating the exact residual
\[
 F-x=\frac{w(b-xk)+ew^2-xH}{D}.
\]
The first two terms cost \(O(\epsilon_d)\); the near \(xH/D\) term
costs \(O(\epsilon_d\sigma L_\ell)\), and the far term
\(O(\epsilon_d)\). Include the original initial value
\(\ell_0H(\sigma)\) to obtain \eqref{deg:landing-direct-value}.

Apply Lemma~\ref{deg:lem-spectral-exterior} to this same intrinsic
landing graph, with starting abscissa \(\sigma\), terminal cut \(x_c\),
and identity observer. Its constant coefficients are
\[
 t=-r_s,\qquad b_{\rm k}=2\gamma_s,\qquad v_{\rm k}=\gamma_s,
 \qquad
 (\mu_2,\mu_1,\mu_0)
   =(3\gamma_s,\epsilon_d\omega,\epsilon_d\kappa B),
 \qquad r_s=\epsilon_d\mathcal R .
\]
Here \(b_{\rm k},v_{\rm k}\) denote the lemma's constant coefficients.
The two exact coefficient graphs give
\[
 H(x)=-r_sx^3+3\gamma_sx^2+\epsilon_d\omega x
                         +\epsilon_d\kappa B,
 \qquad b(x)=x(1+2\gamma_s-r_sx),
\]
so its \(M,Q,\zeta\) are precisely \(H,k,e\) in
\eqref{deg:landing-direct-graph}. On the original enlarged range,
\(2\sigma<x_c\), and \eqref{deg:landing-direct-coefficients} gives
\[
 H(x)/x^2\asymp q_\ell\quad(\sigma\le x\le2\sigma),\qquad
 |\gamma_s|+r_s\le C\epsilon_d,\qquad
 \kappa/\sigma\le C/R_0 .
\]
Thus the lemma's mass bounds hold with \(\rho=q_\ell\), its
\(\Theta\le Cq_\ell\), and, since the observer is the identity,
\(\mathcal T\le C\epsilon_d\). Its smallness conditions follow from
one independent upper bound on \(\tau\) and the original coefficient
bounds, uniformly on the compact and reciprocal graphs. The actual
corridor and terminal height are those already constructed above.

Take the lemma's frozen quadratic-label scale equal to one, so that
\(\beta_\ell=-\sigma^2/2\). By
\eqref{deg:landing-direct-exact-variation}, the terminal slope in this
label is exactly \(\Gamma(\sigma)=-w_\sigma(x_c;\sigma)/\sigma\).
With \(E_\sigma=\sigma\partial_\sigma\),
\eqref{deg:spectral-exterior-rows} gives
\[
 \left|E_\sigma\log\Gamma+\lambda_\gamma\right|
 \le C\{q_\ell L_\ell+\epsilon_d(\sigma/x_c)^e\},\qquad
 \lambda_\gamma=\frac{12\gamma_s}{1+2\gamma_s},\quad
 e=2-\lambda_\gamma>3/2 .
\]
The spectral power \(\sigma^{-\lambda_\gamma}\) has its exact exponent
held fixed in this phase derivative. Since
\(|\lambda_\gamma|\le C\epsilon_d\), \(\sigma<1\), and \(e>1\),
division by \(\sigma\) yields
\[
 |\partial_\sigma\log\Gamma|
 \le C\{(q_\ell/\sigma)L_\ell+\epsilon_d\}.
\]
The original action identity
\(\mathcal D=s_0(\mathcal P/r)\partial_\sigma\), with
\(\mathcal P/r\asymp1\), gives the second row of
\eqref{deg:landing-direct-first-row}. For \(r\ge R_0\),
\(q_\ell=\tau/(\sqrt B\,r)\) and \(s_0\epsilon_d=O(\kappa^2\tau)\)
give its supremum \(C\tau(1+|\log\tau|)\), independently of the later
\(R\) and of the reciprocal parameter.
\end{proof}

\begin{corollary}\label{deg:cor-growing-landing-fixed-product}
On a fixed original band \(0<\sigma_1\le\sigma\le\sigma_2\le\bar\sigma\),
\begin{equation}\label{deg:landing-direct-fixed-product}
 \left\|W_\theta(\sigma/s_0)-\frac{u_c^2-\sigma^2}{2}\right\|_
                  {C^2_\sigma[\sigma_1,\sigma_2]}
 \le C\epsilon_d(1+|\log\epsilon_d|)\longrightarrow0
 \quad\text{as }\epsilon_d\to0 .
\end{equation}
If, independently, the actual scalar inverse phase
\(\Psi_\theta\to\Psi_0\) in \(C^2\) on a fixed buffered interval
containing this band's height image, then
\(\Psi_\theta(W_\theta(\sigma/s_0))\) converges in \(C^2_\sigma\)
to \(\Psi_0((u_c^2-\sigma^2)/2)\), under the hypotheses of that
independent convergence.
\end{corollary}
\begin{proof}
On the fixed band, \eqref{deg:landing-direct-gamma-unit} and
\eqref{deg:landing-direct-first-row} give
\(\Gamma-1,\Gamma_\sigma=O(\epsilon_d(1+|\log\epsilon_d|))\).
Apply these to \(W_\sigma=-\sigma\Gamma\), then use the value
row \eqref{deg:landing-direct-value}. The composition claim is the
ordinary \(C^2\) chain rule on the assumed buffered image; for the
inverse phase use Proposition~\ref{fx:prop-inverse-phase} under its
own limiting hypotheses.
\end{proof}

In coordinates scaled by the starting lip, the fixed terminal section
has a moving coordinate. Its contribution remains in the derivative: if \(x=\sigma t\), \(w=\sigma^2v\), then
\[
 \frac{d}{d\sigma}W_\theta(\sigma/s_0)
 =2\sigma v(x_c/\sigma;\sigma)
  +\sigma^2v_\sigma(x_c/\sigma;\sigma)
  -x_cv_t(x_c/\sigma;\sigma).
\]
The limiting heights lie inside the original buffered
\(\mathcal I_{\rm out}\), including \(u_c^2/2\); the value estimate
places the whole enlarged landing image there after a common shrink.
With the actual raw map \eqref{deg:fixed-exit-actual-output},
\[
 T(r)=\Xi(W_\theta(r)),\qquad
 T'(r)=-s_0^2r\Gamma_\theta(r)\Xi'(W_\theta(r)).
\]
Its raw endpoint, exact power \(\Omega=r_s^{2-\lambda_\gamma}\),
and ordinary \(C^3\) normalized observer bounds are retained.
No coefficient derivative of a moving raw endpoint is asserted.
For the positive unit \(G=\Gamma_\theta\Xi'/(\Omega U_{\rm out})\),
the original chain rule gives the explicit bound
\begin{equation}\label{deg:landing-direct-raw-first-row}
 |\mathcal D\log G|
 \le C\left\{\frac{q_\ell}{r}L_\ell+s_0\epsilon_d
                           +s_0^2\mathcal P(r)\right\}.
\end{equation}
The last term comes from the bounded logarithmic derivative
\(\Xi''/\Xi'\) of the independently constructed output map;
the exact common tail power cancels.
The error in the \(\Gamma_\theta\) derivative tends to zero with
$\tau$; the output-composition error also requires $s_0\to0$.
We now construct the matching independently of this landing estimate.

\subsection{Separate source heights in the growing range}

The positive and negative half-orbits reach the common centre with
exponentially small heights. We construct these two height maps
separately, retaining their linear multipliers and initial polynomial
factors. Equality of their logarithms then gives the matched action
derivative while allowing both centre multipliers to tend to zero.

\begin{lemma}[Actual source matching at degree loss]
\label{deg:lem-growing-matching-direct}
Fix the original compact or reciprocal coefficient ranges, sections,
passive coefficient neighbourhood and buffered positive weak chord
of \eqref{deg:fixed-exit-exact-weak-coordinate}.
For one fixed \(R_0>r_c\), the separate centre-height maps on
\(0<P<p_0\) and \(R_0\le r\le\bar\sigma/s_0\) have the
pure rows \eqref{deg:source-growing-negative-source-rows} and
\eqref{deg:source-growing-positive-source-rows}.
Their product remainder has the total-order-two rows
\eqref{deg:source-growing-product-phase-rows} before matching.
There is a fixed \(R=r_{\rm ge}^{\rm out}>R_0\), chosen before
the independent coefficient bounds, such that every negative lip
with \(R\le r\le\bar\sigma/s_0\) matches one actual input in
the original strict positive chord. On this entire enlarged graph,
\[
 -X\partial_Xr=e_s\frac{\mathcal P(r)}r(1+\zeta),\qquad
 |\zeta|+|X\partial_X\zeta|\le C\tau,
 \qquad e_s=\tau/\chi\asymp1.
\]
All phase derivatives freeze the coefficients and retain the original
centre, action constants, physical scales and raw section windows.
\end{lemma}
\begin{proof}
\emph{The two centre-height maps.}
Use the fixed coefficient field and positive weak chord, with scales
\begin{equation}
 \epsilon=\kappa\tau,\qquad s_0=\kappa\sqrt B,\qquad
 a=\kappa A_d,\qquad x_* =A_d/\sqrt B,\qquad
 r_s=\epsilon\mathcal R,\qquad \gamma_s=\epsilon\mathcal R\nu .
 \label{deg:source-growing-phase-scales}
\end{equation}
Distinguish the positive unit $B$ from $\mathscr B=B_\tau(-r)$.
For the original inner cut $r_c>2$ and fixed $R_0>r_c$, use
\begin{equation}
 R\le r\le\bar\sigma/s_0,\qquad R\ge R_0,\qquad
 0<\bar\sigma<1/2 .
 \label{deg:source-growing-separate-domains}
\end{equation}
Constants are independent of the later $R$. The positive domain is
$0\le P\le p_0$ with its original larger chord buffer.

\smallskip
\noindent\emph{The negative source equation.}

Write \(x=-u=s_0r\), \(w=s_0^2\tau y\), and
\begin{equation}
 P_-(r)=P_s(-r),\qquad K_-(r)=K(-s_0r),\qquad
 e_-(r)=-E(-s_0r).
 \end{equation}
The exact polynomial is
\begin{equation}
 P_-(r)=\frac{\mathcal R}{\sqrt B}
       (1-s_0r)(x_*+r)(c+s_0r)
 \quad\hbox{on the compact graph},
 \label{deg:source-growing-negative-polynomial-compact}
\end{equation}
and
\begin{equation}
 P_-(r)=\frac d{\sqrt B}
       (1-s_0r)(x_*+r)(1+js_0r)
 \quad\hbox{on the reciprocal graph}.
 \label{deg:source-growing-negative-polynomial-remote}
\end{equation}
The identities $\mathcal RA_dc=B$ and $B=dA_d$, respectively,
give exactly the original $P_s(-r)$.
Set
\begin{equation}
 \mathscr B(r)=1+2\gamma_s-\epsilon\mathcal R s_0r,\quad
 \mathfrak d(r)=\frac{P_-(r)}{r\mathscr B(r)},\quad
 \mathcal D=\mathfrak d(r)\partial_r,\quad
 h_1(r)=\frac{P_-'(r)}{r\mathscr B(r)} .
 \label{deg:source-growing-negative-action-operator}
\end{equation}
Use the chosen coefficient-only real action offset \(A_{c,0}(\theta)\)
and define
\begin{equation}
 \mathcal A_-(r)=\int_{r_c}^r\frac{dt}{\mathfrak d(t)}
                       -A_{c,0},\qquad \mathcal D\mathcal A_-=1.
 \label{deg:source-growing-chosen-action-offset}
\end{equation}
Direct differentiation of \eqref{deg:source-growing-negative-polynomial-compact}, \eqref{deg:source-growing-negative-polynomial-remote}, and \eqref{deg:source-growing-negative-action-operator} gives
\begin{equation}
 \begin{gathered}
 P_-\asymp1+r,\quad \mathscr B\asymp1,\quad
 \mathfrak d\asymp1,\quad
 |\mathcal D^m\mathfrak d|\le C_m,\quad
 |\mathcal D^mh_1|\le C_m/r\quad(0\le m\le3),\\
 |K_--1|\le C\epsilon,\quad
 |\mathcal D^mK_-|\le C_m\tau\quad(1\le m\le3).
 \end{gathered}
 \label{deg:source-growing-negative-coefficient-bounds}
\end{equation}
For the polynomial monotonicity used below, put \(\sigma=s_0r\).
On the compact graph,
\[
 P_-'=\frac{\mathcal R}{\sqrt B}
 \{c(1-2\sigma)+2\sigma-3\sigma^2
                         +a(1-c-2\sigma)\},
\]
whereas on the reciprocal graph
\[
 P_-'=\frac d{\sqrt B}
 \{1-2\sigma+j\sigma(2-3\sigma)
                         +a(j(1-2\sigma)-1)\}.
\]
The fixed bound \(\bar\sigma<1/2\) and a common \(a\)-shrink
give \(P_-'\ge c\) in either graph.  Since
\(\mathscr B'=-\epsilon\mathcal R s_0\le0\), it follows that
\begin{equation}\label{deg:degree-negative-polynomial}
 \mathcal P:=P_-/\mathscr B\asymp1+r,\qquad
 \mathcal P'\ge c .
\end{equation}

Indeed, $P_-'$, $P_-''$, $P_-'''$ are bounded by
$C,Cs_0,Cs_0^2$, with $s_0\le\bar\sigma/r$. Also
\begin{equation}
 K_-=1-A(\gamma_s-\epsilon\mathcal R s_0r)
           +C(\gamma_s-\epsilon\mathcal R s_0r)^2.
 \label{deg:source-growing-negative-K}
\end{equation}
Its positive-order action rows are $O(\epsilon s_0)=O(\tau)$.
On the reciprocal graph, $\mathcal R=dj$ and
$\gamma_s=\epsilon d(j-ja-1)/3$ are regular at $j=0$.

From the original rational graph equation one obtains exactly
\begin{equation}
 y_r=\frac{y\{r\mathscr B+e_-s_0\tau y\}}
                 {\tau(P_-+K_-y)}.
 \label{deg:source-growing-negative-scalar-equation}
\end{equation}
Along the inverse passage introduce accumulated action \(q\)
by \(r_q=-\tau\mathfrak d(r)\), starting at \(r\), and put
\(W=y/P_-\).  The source lip is exactly \(W(0)=\ell_0\).
Then
\begin{equation}
 (\log W)_q=
 -\frac{1+\varpi(r)W}{1+K_-(r)W}+\tau h_1(r)
 =-1+\tau h_1(r)+G(r,W),
 \label{deg:source-growing-stable-ratio-equation}
\end{equation}
where
\begin{equation}
 \varpi=e_-s_0\tau\mathfrak d,\qquad
 G(r,W)=\frac{(K_--\varpi)W}{1+K_-W}.
 \label{deg:source-growing-stable-ratio-nonlinearity}
\end{equation}
Here \(|e_-|\le C\epsilon^2\); direct differentiation gives
\begin{equation}
 |\mathcal D^m\varpi|\le C_m\tau,\qquad 0\le m\le3.
 \label{deg:source-growing-stable-coefficient-rows}
\end{equation}
These derivative bounds apply to the exact nonlinear coefficient.

\smallskip
\noindent\emph{Negative source variations.}\par\nopagebreak[4]

The endpoint at the original inner cut is
\begin{equation}
 T(r)=\frac1\tau\int_{r_c}^r\frac{dt}{\mathfrak d(t)}
     =\frac{\mathcal A_-(r)+A_{c,0}}{\tau}.
 \label{deg:source-growing-inward-terminal}
\end{equation}
It satisfies
\begin{equation}
 T\ge c(r-r_c)/\tau,\qquad
 \mathcal DT=1/\tau,\qquad
 \mathcal D^2T=\mathcal D^3T=0 .
 \label{deg:source-growing-inward-terminal-action-rows}
\end{equation}
Apply Lemma~\ref{deg:lem-stable-ratio-ledger} with
\((\kappa,H,Q,\omega,\rho)=(\tau,h_1,K_-,\varpi,\tau)\)
and density \(L=G\). Its phase derivative is the initial action
derivative \(\delta\) at fixed \(q\):
\begin{equation}\label{deg:source-growing-initial-action-flow}
 \delta f(r(q;r))=(\mathcal Df)(r(q;r)),\qquad
 \partial_q f=-\tau\mathcal Df,
\end{equation}
because \(\mathcal A_-(r(q;r))=\mathcal A_-(r)-\tau q\).
In particular \(\delta\) commutes with \(\partial_q\).
The polynomial factors remain positive on the whole route
\(r_c\le v\le r\), with \(s_0v\le\bar\sigma<1/2\), and on
a fixed enlargement at \(r_c\). Equations
\eqref{deg:source-growing-negative-coefficient-bounds} and
\eqref{deg:source-growing-stable-coefficient-rows} give the required
phase and mixed rows with \(\alpha=0\). Since \(K_-\) is close
to one and \(\varpi,\tau h_1\) are small, both
\(\mathcal F/W\) and \(\mathcal F_W\) have a fixed negative
margin on \(0\le W\le\ell_0\). The lemma therefore gives
\begin{equation}\label{deg:source-growing-stable-ratio-corridor}
 0<W(q;r)\le\ell_0e^{-cq},\qquad 0\le q\le T(r),
\end{equation}
and the fixed-clock rows
\begin{equation}
 |\delta^iG|\le C\tau(1+q)^Ne^{-cq}\quad(i=1,2),
 \qquad |G|+|G_q|\le C(1+q)^Ne^{-cq}.
 \label{deg:source-growing-stable-kernel-rows}
\end{equation}
The inverse passage thus reaches the inner cut positively, with positive
denominator, and its lip is the first ratio hit.
Define on this actual passage
\begin{equation}
 J(r)=\int_0^{T(r)}G(r(q;r),W(q;r))\,dq .
 \label{deg:source-growing-nonlinear-action-integral}
\end{equation}
The ratio equation gives exactly
\(G=-K_-W_q+(\tau K_-h_1-\varpi)W\). Integration by parts
therefore isolates the nonlinear mass:
\begin{equation}\label{deg:source-growing-exact-action-integral}
 J=\ell_0K_-(r)-K_-(r_c)W(T;r)+\int_0^T QW\,dq,
 \qquad Q=\tau(K_-h_1-\mathcal DK_-)-\varpi .
\end{equation}
Here \(\mathcal D^iQ=O(\tau)\), \(0\le i\le2\), using
\(\mathcal D^3K_-\) in the second row. The same fixed-clock
estimates give an integrable \(C\tau(1+q)^Ne^{-cq}\) bound for
\(QW\) and its phase rows through order two. Apply the endpoint formulas
\eqref{deg:stable-ratio-terminal-rows} to this density and the
terminal chain rules of Lemma~\ref{deg:lem-stable-ratio-ledger}
to \(K_-(r_c)W(T;r)\), with
\(T_1=1/\tau\), \(T_2=0\). The fixed gap \(R_0-r_c>0\)
makes all endpoint terms \(\tau^{-N}(1+T)^Ne^{-cT}\) flat.
Together with \(K_-(r)-1=O(\epsilon)=O(\tau)\), this proves
\begin{equation}
 |J-\ell_0|+|\mathcal DJ|+|\mathcal D^2J|\le C\tau .
 \label{deg:source-growing-nonlinear-action-rows}
\end{equation}
The constants depend on \(R_0-r_c>0\), not on the later \(R\).

The polynomial-memory identity
\eqref{deg:stable-ratio-polynomial-memory} retains the initial factor:
\begin{equation}
 \int_0^T\tau h_1(r(q;r))\,dq
             =\log P_-(r)-\log P_-(r_c).
 \label{deg:source-growing-polynomial-memory}
\end{equation}
Consequently the actual height at the inner cut is
\begin{equation}
 y_c(r)=\ell_0P_-(r)e^{-T(r)}e^{J(r)} .
 \label{deg:source-growing-inner-cut-height}
\end{equation}
Uniformly for \(r\ge R_0\),
\begin{equation}
 0<y_c(r)\le C(1+r)e^{-c(r-r_c)/\tau}
           \le C_N\tau^N\quad\hbox{for every fixed }N .
 \label{deg:source-growing-inner-cut-flatness}
\end{equation}
Its first two logarithmic rows are
\begin{equation}
 \mathcal D\log y_c=h_1-1/\tau+\mathcal DJ,\qquad
 \mathcal D^2\log y_c=\mathcal Dh_1+\mathcal D^2J .
 \label{deg:source-growing-inner-cut-log-rows}
\end{equation}

We now continue to the original common centre, without
renormalizing it to an image-dependent section.
On \(0\le v\le r_c\), \(P_-(v)\) and \(\mathscr B(v)\)
are positive units.  The exact linear multiplier is
\begin{equation}
 L_-(v)=\exp\left\{-\frac1\tau
               \int_v^{r_c}\frac{t\mathscr B(t)}{P_-(t)}\,dt
                  \right\}\le1,\qquad
 A_{\rm core,-}=\int_0^{r_c}\frac{t\mathscr B(t)}{P_-(t)}\,dt.
 \label{deg:source-growing-negative-core-linear-factor}
\end{equation}
The exact equation on this fixed compact segment is
\(y_v=A_-y+\mathcal K_-y^2\), where
\begin{equation}
 A_- =\frac{v\mathscr B}{\tau P_-},\qquad
 \mathcal K_-(v,y)=\frac{e_-s_0}{P_-+K_-y}
       -\frac{v\mathscr B K_-}{\tau P_-(P_-+K_-y)} .
 \label{deg:source-growing-negative-core-equation}
\end{equation}
On the original fixed small height buffer, all normal derivatives
through order three of \(\mathcal K_-\) are bounded by \(C/\tau\).
Its inward linear propagator is exactly \(L_-\le1\), and
\eqref{deg:source-growing-inner-cut-flatness} gives \(y_c/\tau\ll1\).
Lemma~\ref{deg:lem-compact-height-continuation}, with
\(\kappa_c=\tau\), gives
an actual positive centre hit and
\(y(0;y_c)=L_-(0)y_c\mathcal U_-(y_c)\), where
\begin{equation}
 |\log\mathcal U_-|+
 \sum_{i=1}^2|(y_c\partial_{y_c})^i\log\mathcal U_-|
                                  \le Cy_c/\tau .
 \label{deg:source-growing-negative-core-height-rows}
\end{equation}
This uses the full multiplier, with no reciprocal bound on it, and
holds for either sign of \(e_-(0)\).

Use the same centre-height scale as on the positive side:
\begin{equation}
 Z_-(r)=\frac{w(0;r)}{s_0^2\tau}
 =\ell_0P_-(r)
       \exp\left\{-\frac{\mathcal A_-(r)+A_{c,0}
                             +A_{\rm core,-}}{\tau}\right\}
       e^{J(r)}\mathcal U_-(y_c(r)).
 \label{deg:source-growing-negative-centre-height}
\end{equation}
Choose the coefficient-only logarithmic constant
\begin{equation}
 c_-=\log\ell_0-(A_{c,0}+A_{\rm core,-})/\tau .
 \label{deg:source-growing-negative-coefficient-constant}
\end{equation}
The separate negative remainder is
\begin{equation}
 S_-(r)=\log Z_-(r)+\mathcal A_-(r)/\tau
                            -\log P_-(r)-c_-
       =J(r)+\log \mathcal U_-(y_c(r)).
 \label{deg:source-growing-negative-source-remainder}
\end{equation}
The chain rule, \eqref{deg:source-growing-inner-cut-flatness}--\eqref{deg:source-growing-inner-cut-log-rows}, and \eqref{deg:source-growing-negative-core-height-rows} bound the two
core contributions by \(Cy_c/\tau^2\) and \(Cy_c/\tau^3\).
Combining with \eqref{deg:source-growing-nonlinear-action-rows} gives
\begin{equation}
 |S_-|\le C,\qquad
 |\mathcal DS_-|+|\mathcal D^2S_-|\le C\tau .
 \label{deg:source-growing-negative-source-rows}
\end{equation}
The physical centre height is therefore
\begin{equation}\label{deg:source-growing-physical-negative-height}
 z_-(r)=r_s^2s_0^2\tau e^{c_-}P_-(r)
                  e^{-\mathcal A_-(r)/\tau}e^{S_-(r)}.
\end{equation}
The factor \(e^{S_-}\), its inverse and its first two action
rows are uniformly bounded.


\smallskip
\noindent\emph{Positive source variations.}

The positive weak coordinate from
\eqref{deg:fixed-exit-exact-weak-coordinate} satisfies
\begin{equation}
 \Phi(P)=\{z_+y(P)\}^{\chi},\qquad
 \chi=\tau X_d,\qquad z_+^{\pm1}+X_d^{\pm1}\le C .
 \label{deg:source-growing-positive-exact-coordinate}
\end{equation}
The local invariant gives
\begin{equation}
 0\le y(P)\le C\vartheta^{1/\chi},\quad
 0<\vartheta<1,\quad 0\le P\le p_0,
 \label{deg:source-growing-positive-weak-flatness}
\end{equation}
and \(\Phi'>c>0\) with its stated uniform phase bounds.
Fix \(P_\dagger\in(0,p_0)\), independently of coefficients,
and set \(X=\Phi(P)/\Phi(P_\dagger)\).  On \(P>0\),
\begin{equation}
 \frac{y(P)}{y(P_\dagger)}=X^{1/\chi},\qquad
 \mathcal E:=X\partial_X=\chi^{-1}y\partial_y .
 \label{deg:source-growing-positive-log-phase}
\end{equation}
Use this exact logarithmic derivative.

Construct \(Z_+(P)\) independently by following positive time
from \(u_+=a(1-\eta_+)\), with \(w=a^2y(P)\), to \(u=0\).
In \(v=u/a\), \(v_+=1-\eta_+\), put
\begin{equation}
 h(av)/a^2=\tau H_+(v),\qquad
 L(av)/a=v\mathscr B_+(v),\qquad K_+(v)=K(av).
 \end{equation}
The original exact functions are
\begin{equation}
 H_+(v)=\frac{\mathcal R}{A_d}
        (1+av)(1-v)(c-av)
 \end{equation}
on the compact graph, and
\begin{equation}
 H_+(v)=\frac d{A_d}(1+av)(1-v)(1-jav)
 \end{equation}
on the reciprocal graph.  On \(0\le v\le v_+\),
\(H_+,\mathscr B_+,K_+\) are positive units, with uniform
finite coefficient-function bounds.  In particular the
positive core has no inverse \(\kappa\) scale.
Its exact scalar equation is
\begin{equation}
 y_v=\frac{y\{v\mathscr B_++aE(av)y\}}
                         {\tau H_++K_+y}.
 \label{deg:source-growing-positive-core-equation}
\end{equation}
Its linear multiplier and positive core action are
\begin{equation}
 L_+(v)=\exp\left\{-\frac1\tau
              \int_v^{v_+}\frac{t\mathscr B_+(t)}{H_+(t)}\,dt
                    \right\}\le1,\qquad
 A_{\rm core,+}=\int_0^{v_+}\frac{t\mathscr B_+(t)}{H_+(t)}\,dt.
 \label{deg:source-growing-positive-core-linear-factor}
\end{equation}
Freeze the coefficients and set \(y(v)=\tau^2\eta(v)\),
with initial height \(\eta_{\rm in}=y(P)/\tau^2\) at the same cut
\(v_+\). Exact division gives
\begin{equation}\label{deg:source-growing-positive-core-rescaled-equation}
 \eta_v=\frac{v\mathscr B_+}{\tau H_+}\eta
                         +\widetilde{\mathcal K}_+(v,\eta)\eta^2,
 \qquad
 \widetilde{\mathcal K}_+
 =\frac{\tau aE(av)}{H_++\tau K_+\eta}
 -\frac{v\mathscr B_+K_+}{H_+(H_++\tau K_+\eta)}.
\end{equation}
Choose a fixed small buffer \(0\le\eta\le\eta_b<\inf X_d^2\).
The nonlinear coefficient and its first three ordinary height
derivatives are bounded uniformly. Indeed the
displayed denominators are positive units and each such derivative
introduces a bounded factor \(\tau K_+\).
Then \(\tau^2\eta_b<\chi^2\) lies inside the original height range;
\eqref{deg:source-growing-positive-weak-flatness} puts every original
\(0<P\le p_0\) in a smaller initial subbuffer after a common
\(\tau\)-shrink.
The inward propagator remains exactly \(L_+\le1\), without a
reciprocal reserve. Thus
Lemma~\ref{deg:lem-compact-height-continuation}, with
\(\kappa_c=1\), supplies the actual positive centre hit
and its Euler-height rows. Since
\(\eta_{\rm in}\partial_{\eta_{\rm in}}=y_+\partial_{y_+}\)
for \(y_+=y(P)\), these give
\begin{equation}
 |\log \mathcal U_+(y_+)|+
 \sum_{i=1}^2|(y_+\partial_{y_+})^i\log \mathcal U_+(y_+)|
                                      \le Cy_+/\tau^2 ,
 \label{deg:source-growing-positive-core-height-rows}
\end{equation}
where \(\mathcal U_+\) is the normalized factor at \(v=0\).
The actual central height is
\begin{equation}
 Z_+(P)=C_+(\theta)y(P)\mathcal U_+(y(P)),\qquad
 C_+(\theta)=\frac{a^2}{s_0^2\tau}
                                      e^{-A_{\rm core,+}/\tau}.
 \label{deg:source-growing-positive-centre-height}
\end{equation}
This retains exactly \(a^2/(s_0^2\tau)=A_d^2/(B\tau)\).
The prefactor of the exponential is a coefficient-only unit times
\(\tau^{-1}\).
Scalar first variation gives \(Z_+'(P)>0\) for \(P>0\).

Define on the original independent positive domain
\begin{equation}
 S_+(P)=\log Z_+(P)-\log Z_+(P_\dagger)-\chi^{-1}\log X
 =\log \mathcal U_+(y(P))-\log \mathcal U_+(y(P_\dagger)).
 \label{deg:source-growing-positive-source-remainder}
\end{equation}
The coefficient-only multiplier cancels as an exact reference
difference.  Equations \eqref{deg:source-growing-positive-log-phase} and \eqref{deg:source-growing-positive-core-height-rows} give
\begin{equation}
 |\mathcal E^iS_+|
 \le C\chi^{-i}\tau^{-2}y(P)
 \le C\tau^{-i-2}\vartheta^{1/\chi},\qquad i=1,2 .
 \label{deg:source-growing-positive-source-rows}
\end{equation}
The rows vanish continuously at $P=0$. This positive map remains
independent of the negative image.

\smallskip
\noindent\emph{Admission to the positive input chord.}

Choose any fixed \(p_*\in(0,p_0)\) inside the original counted
weak chord, with its larger original raw chord buffer.
Put \(j(P)=\Phi(P)/\Phi'(P)\).  Its uniform phase bounds give
\(j'(P)=1+O(P)\); choose \(p_*\) small enough that \(j'(P)>3/4\)
on \(0\le P\le p_*\).  The endpoint is fixed before \(R\).
By \(\Phi'>c\), \(\Phi(p_*)\) is bounded below by a fixed
positive constant.  Equations \eqref{deg:source-growing-positive-exact-coordinate}, \eqref{deg:source-growing-positive-core-linear-factor}, and \eqref{deg:source-growing-positive-centre-height} give
\begin{equation}
 -\tau\log Z_+(p_*)\le C_*(p_*) ,
 \label{deg:source-growing-positive-endpoint-action-budget}
\end{equation}
uniformly in \(\kappa,\tau,j\) after a common small bound.
Indeed, $-\tau\log y(p_*)=-(\tau/\chi)\log\Phi(p_*)
+\tau\log z_+$ and $A_{\rm core,+}$ are bounded; the remaining
scale costs $O(\tau(1+|\log\tau|))$.

On the negative side, \eqref{deg:source-growing-negative-centre-height} gives
\begin{equation}
 -\tau\log Z_-(R)\ge
 c(R-r_c)-\tau\log\{\ell_0P_-(R)\}-C\tau .
 \label{deg:source-growing-negative-endpoint-action-budget}
\end{equation}
First choose a fixed \(R=r_{\rm ge}^{\rm out}>R_0\), larger than any
already prescribed original growing threshold, so that
\(c(R-r_c)>C_*(p_*)+2d_*\) for a fixed \(d_*>0\).
Only afterwards decrease the coefficient bounds so that
\(\tau\log(1+R)\) and the fixed terms in \eqref{deg:source-growing-negative-endpoint-action-budget} are below
\(d_*\).  This preserves the full independent coefficient
family.
Choose the counted threshold \(r_{\rm ge}>R\) with a fixed
positive buffer before decreasing the coefficients. It follows that
\begin{equation}
 0<Z_-(R)<e^{-d_*/\tau}Z_+(p_*).
 \label{deg:source-growing-positive-image-inclusion}
\end{equation}
Furthermore \eqref{deg:source-growing-negative-source-remainder}--\eqref{deg:source-growing-negative-source-rows} imply
\begin{equation}
 \mathcal D\log Z_-=-1/\tau+h_1+\mathcal DS_-<0 .
 \label{deg:source-growing-negative-centre-order}
\end{equation}
Thus every actual growing lip in \eqref{deg:source-growing-separate-domains} has centre height
strictly between zero and \(Z_+(p_*)\).
The continuous strictly increasing positive map \(Z_+\)
maps \([0,p_*]\) onto \([0,Z_+(p_*)]\).
There is a unique \(P(r)\in(0,p_*)\) with
\(Z_+(P(r))=Z_-(r)\).
Scalar uniqueness at the common centre identifies the two
actual graph segments with the same orbit.

The local weak chord is the one parameterized by the original
raw input
\begin{equation}
 P=(x_{\rm raw}-x_b)/s_{\rm in},\qquad
 s_{\rm in}=\Omega a^{2-\lambda_\gamma}U_{\rm in},\qquad
 \Omega=r_s^{2-\lambda_\gamma},\quad U_{\rm in}^{\pm1}\le C.
 \label{deg:source-growing-positive-raw-input-scale}
\end{equation}
The independently constructed stopped outer chord and its
local invariant give the actual continuation to this raw
section for all \(0\le P\le p_0\), with its original buffers.
The matched inputs lie in the strict interior \(P<p_*\).
Both local points are on the same invariant level in the
original coordinate box; the monotone positive invariant
coordinate gives the intervening orbit segment there.
Thus every lip has the actual positive-time continuation from raw
input through the positive middle cut and centre to the negative lip.
The separate annular overlap remains the earlier construction.

\smallskip
\noindent\emph{The matched action derivative.}

On the separate original product domain define
\begin{equation}
 R_{\rm prod}(P,r)=\log P_-(r)+S_-(r)-S_+(P).
 \label{deg:source-growing-product-remainder}
\end{equation}
It remains defined whether or not the two images intersect.
Its phase rows satisfy
\begin{equation}
 \begin{gathered}
 |\mathcal ER_{\rm prod}|+|\mathcal E^2R_{\rm prod}|
                  \le C\tau^{-4}\vartheta^{1/\chi},\\
 |\mathcal DR_{\rm prod}|+|\mathcal D^2R_{\rm prod}|\le C,
 \qquad \mathcal E\mathcal DR_{\rm prod}=0 .
 \end{gathered}
 \label{deg:source-growing-product-phase-rows}
\end{equation}
Although \(\log P_-(r)\) can grow, its two action derivatives have
the bounds needed here.
On the actual matched graph the exact identity is
\begin{equation}
 \mathcal A_-(r)=
 \tau\{c_--\log Z_+(P_\dagger)\}
             -e_s\log X+\tau R_{\rm prod}(P,r),
 \qquad e_s=\tau/\chi=1/X_d .
 \label{deg:degree-growing-matching}
\end{equation}
The constant retains all coefficient-only terms; $e_s$ is a positive unit.

Define separate product functions
\begin{equation}
 N_-=1-\tau\{h_1+\mathcal DS_-\},\qquad
 N_+=1+\chi\mathcal ES_+,\qquad n(P,r)=N_+(P)/N_-(r).
 \label{deg:source-growing-separate-action-numerators}
\end{equation}
For a common small \(\tau\) bound,
\begin{equation}
 N_\pm\ge1/2,\quad
 |N_--1|+|\mathcal DN_-|\le C\tau,\quad
 |N_+-1|+|\mathcal EN_+|
                       \le C\tau^{-3}\vartheta^{1/\chi}.
 \label{deg:source-growing-action-numerator-rows}
\end{equation}
Differentiate equality of the two actual centre heights:
\(\mathcal E\log Z_+=N_+/\chi\) and
\(\mathcal D\log Z_-=-N_-/\tau\).
It gives exactly
\begin{equation}
 -\mathcal E r=e_s\mathfrak d(r)n(P,r)
     =e_s\frac{\mathcal P(r)}r(1+\zeta),\qquad \zeta=n-1 .
 \label{deg:source-growing-actual-matching-slope}
\end{equation}
The matched graph is consequently strictly decreasing. To differentiate
the quotient along this graph, both separate variables must be varied:
\begin{equation}
 \mathcal E_{\rm graph}\zeta=
 \mathcal E_+n-e_sn\mathcal D_-n .
 \label{deg:source-growing-actual-quotient-first-row}
\end{equation}
Here \(\mathcal E_+\) and \(\mathcal D_-\) act on the separate
positive and negative variables of the product.
The second pure derivatives in \eqref{deg:source-growing-negative-source-rows}
and \eqref{deg:source-growing-positive-source-rows} bound precisely
the two derivatives in \eqref{deg:source-growing-actual-quotient-first-row}.  Thus
\begin{equation}
 |\zeta|+|\mathcal E_{\rm graph}\zeta|
 \le C\{\tau+\tau^{-3}\vartheta^{1/\chi}\}
 \le C\tau .
 \label{deg:source-growing-actual-quotient-estimates}
\end{equation}
This is uniform on the whole enlarged graph with independent
$\kappa,\tau,j$; the error tends to zero with $\tau$.

The normalized source heights retain the physical conversion
\(z(\gamma_s)=r_s^2s_0^2\tau Z\), the exact raw power
\(\Omega\), the positive and negative tail quartets on the
reciprocal graph
\begin{equation}
 \left(-\frac{\tau d}{A_d},
       \frac{\tau d(1-a+ja)}{A_d},\gamma_s,r_sa\right),
 \qquad
 \left(-\kappa^2\tau dA_d,
       \kappa\tau d(1-a+ja),\gamma_s,r_s\right),
 \label{deg:source-growing-remote-raw-quartets}
\end{equation}
and the passive coefficient restrictions used to construct the raw
chords.  No estimate of \(\gamma_s\log r_s\), no derivative
with respect to a coefficient, and no frozen moving raw endpoint
is used.  At \(j=0\) only the one-sided normalized estimates
are intended; actual raw sections are asserted for \(j>0\).
\end{proof}


\begingroup
\let\degreecarriersubsection\subsection
\let\degreecarriersubsubsection\subsubsection
\let\section\degreecarriersubsection
\let\subsection\degreecarriersubsubsection
\section{One physical interval and four passage regions}

The preceding constructions use different phase coordinates. We now show
that their domains cover one connected interval of actual through
trajectories, with open overlaps. Their derivative estimates can then
be combined for the same affine displacement.

\begin{lemma}[A common entrance and gate choice]
\label{p2:lem-joint-entrance}
One auxiliary choice for the four passage regions, followed by
independent small coefficient bounds, has the following properties.
The incoming raw height $h$ is $x_{\rm raw}$ in
\eqref{deg:fixed-exit-natural-input} and $x_{\rm in}$ in
Lemma~\ref{lem:actual-dulac-matching}; its stopped value is $x_b$.
Thus the canonical weak input is exactly
\begin{equation}\label{p2:joint-raw-input}
 P=q(h)=\frac{h-x_b}{s_{\rm in}},\qquad
 s_{\rm in}=\Omega a^{2-\lambda_\gamma}U_{\rm in}>0,
 \qquad \Omega=r_s^{2-\lambda_\gamma},\quad U_{\rm in}^{\pm1}\le C.
\end{equation}
This is the affine coordinate on the original raw section, distinct
from the displacement coordinate $x=h-a_\theta$; no new input
observer is inserted.

Let $\widehat p_0=p_0$ be a fixed buffered raw cutoff in
\eqref{deg:fixed-exit-middle-invariant}, and choose
$0<\widehat p_*<\widehat p_0$ so that
$\mathfrak j'=(\Phi_\theta/\Phi_\theta')'>3/4$ on
$[0,\widehat p_*]$. Set $q_{\rm ref}=\widehat p_*/4$ and use
$P_\dagger=q_{\rm ref}$ in
\eqref{deg:source-growing-positive-log-phase}. The final negative
thresholds are denoted by
$\widehat r_{\rm ge}^{\rm out}=r_{\rm ge}^{\rm out}$ and
$\widehat r_{\rm ge}=r_{\rm ge}$, with a fixed positive gap. Every
actual first-lip input with $r_\ell\ge r_{\rm ge}$ satisfies
\begin{equation}\label{p2:joint-entrance-buffer}
 0<P<q_{\rm ref}/2<q_{\rm ref}<\widehat p_*<\widehat p_0,
\end{equation}
including first hits outside the growing product range.

The source constant $\delta_\sharp(L)$ is selected once for the full
bounded-source domain and relative negative-curvature estimate. Its gates
are those below in Lemma~\ref{p2:lem-bounded-source-budget}.
The final $r_{\rm ge}$ precedes $R_{\rm BD}$ and its annular
overlap; all fixed gates and phase buffers precede the final coefficient
bounds. No comparison of the independent scales is imposed.
\end{lemma}
\begin{proof}
Fix the original raw windows and a finite $V_1>1$, say $V_1=2$.
Choose the local cuts and ordinary $C^3$ section maps for this value in
Lemma~\ref{lem:actual-dulac-matching}, retaining its enlarged
interval $0\le V\le V_1+1$. Keep these choices throughout.
Use the strict source block and the two comparison orbits of
Proposition~\ref{bd:prop-bounded-action}: fix its derivative constants,
then $B_-$, then $\delta_\sharp(L)$ for both the domain and relative
estimate, and finally $\delta_1,\delta_2$ and the gates below.
Fix $r_c>2$ and choose $\ell_0$ below the ratio of a fixed real
comparison orbit at $-r_c$, as in
\eqref{deg:degree-ratio-monotonicity}.

Fix a source envelope $\bar\sigma_{\rm src}=1/3$ for the separate
centre-map estimates on $r_c\le r\le\bar\sigma_{\rm src}/s_0$ in
Lemma~\ref{deg:lem-growing-matching-direct}; its derivative constants
are fixed before the positive cutoff and negative thresholds below.
The growing and output intervals will be chosen inside this envelope.
The independent stopped outer chord and local invariant give the
positive input map on $[0,p_0]$ in
\eqref{deg:fixed-exit-natural-input}--\eqref{deg:fixed-exit-exact-weak-coordinate}.
Choose the smaller fixed $\widehat p_*$ above so that the positive
action at $P_0=q_{\rm ref}/2$ also exceeds the fixed middle
variational exponent by two. This is possible since
$\Phi_\theta(P)=PU_\theta(P)$ with uniform positive units and
$\tau/\chi$ a positive unit. These choices precede the final negative
threshold.

The full first-hit cutoff following \eqref{deg:degree-phase-order}
provides a fixed $R(P_0)$. Choose $R=r_{\rm ge}^{\rm out}$ larger
than this and every earlier growing threshold, and so that
$c(R-r_c)>C_*(P_0)+2d_*$ for a fixed $d_*>0$ in
\eqref{deg:source-growing-positive-endpoint-action-budget}.
Choose $r_{\rm ge}>R$ with a fixed positive buffer.
The separate centre maps give
\eqref{deg:source-growing-positive-image-inclusion} with endpoint
$P_0$ instead of $p_*$ after independent coefficient shrink.
Their monotonicity and scalar uniqueness therefore construct the
actual product reference $r_\ell=r_{\rm ge}$ with raw input
$0<P_{\rm ge}<P_0$, in the original buffered chord. The local
invariant supplies its positive-time continuation to the raw section,
as in \eqref{deg:source-growing-positive-raw-input-scale}.

On the full first-hit domain of the original through interval,
\eqref{deg:degree-phase-order} gives $r_{\ell,h}<0$.
Thus $r_\ell(h)\ge r_{\rm ge}$ implies
$h\le h_{\rm ge}=x_b+s_{\rm in}P_{\rm ge}$, and the same positive
affine input gives $0<P(h)\le P_{\rm ge}<P_0$.
This proves \eqref{p2:joint-entrance-buffer} throughout the first-hit
domain, including inputs beyond the product range. The extension follows
from order on the full domain; the product derivative estimates remain
restricted to their original range.

Only now choose $R_{\rm BD}$ so the compact strong projection of
$r_\ell=r_{\rm ge}$ lies strictly below $2R_{\rm BD}$.
The transverse inverse and first-hit order give a fixed
$r_{\rm ge}^{\sharp}>r_{\rm ge}$ in the enlarged annular overlap.
The annular proof now fixes its own outer cuts \(L_{\rm a}\), after
this output bound and its limiting-action envelope. The earlier source
cut \(L\), \(h_L\), \(\delta_\sharp(L)\), and literal gates stay fixed;
both passages continue to the same raw ports. The original weak cuts
satisfy \(c_+=p_{\rm a}/8<p_-=p_{\rm a}/2\) and
\(c_-=p_{\rm a}/8<r_{\rm a}/2\), so the annular root-gap construction
uses this same auxiliary tuple. Only final independent coefficient bounds
are reduced to admit the new positive-height tail cuts.
Choose $0<\sigma_1<\sigma_2<\sigma_0<1/4$, the common
negative section $u_c$, its larger
$\sigma_0<\bar\sigma\le\bar\sigma_{\rm src}$ buffer, with
$\bar\sigma<|u_c|/2$, and strict
action gap $\Delta$ as in the construction after
\eqref{deg:degree-phase-order}. After these fixed choices, shrink
the independent coefficients so the fixed annular overlap lies in
the counted product range and $s_0r_{\rm ge}^{\sharp}<\sigma_1$.
All source, raw and shooting buffers retain strict margins.
All these restrictions are independent upper bounds on the coefficients.
After they have been imposed, the finite multiplier will be chosen.
\end{proof}

\begin{proposition}[One actual carrier for the four passage regions]
\label{p2:prop-unified-carrier-cover}
Use the common choice of Lemma~\ref{p2:lem-joint-entrance}, the original field,
through word, sections, raw windows, $L,h_L$ and closing maps.
The source constant $\delta_\sharp(L)$ is the once-selected constant
for the full original bounded-source domain and its relative sign;
$\delta_1,\delta_2$ keep their literal values from that tuple.
Let $I_\theta$ be the connected interval of actual through inputs in
this word, within the original buffered windows and positive graph
denominator domain, with
\[
             0<H_\theta(h):=z(\gamma_s;h)<z_{\rm cap}.
\]
Here $z_{\rm cap}$ is the original common fixed raw-height cap,
chosen within the standing high-height bound. The stopped trace $x_b$
is excluded. Write $h^0$ for the actual trace of $h_L=0$ and
$I_\theta^- = I_\theta\cap\{h\le h^0\}$.

\begin{enumerate}
\item The original through interval $(x_b,h^0)$ has the strictly increasing
shooting coordinate $W_\theta(h)=w(u_c;h)$ with image
$(0,W_\theta(h^0))$, where $W_\theta(h^0)\to u_c^2/2$ uniformly.
Thus every point of the full original enlarged chord
$\overline{\mathcal J_c^{\sharp\sharp}}$ corresponds to an actual through orbit,
with a strict height buffer. Its backward shot and both raw continuations
are that same orbit.
\item Let $h_S$ be the trace of $h_L=-\delta_2$ and let $h_A,h_B$
be the gate traces of $r_b=r_a/2,2R_{\rm BD}$, respectively. Put
\[
 B_\theta=I_\theta^-\cap([h_S,h^0]\cup[h_B,h_A]),\qquad
 h_T=\min\{h_S,h_B\}.
 \tag{PCov1}
\]
Let $I_{\ell,\theta}$ be the original first-ratio-hit domain. On it put
$\sigma(h)=s_0r_\ell(h)$. Wherever $r_\ell\ge r_{\rm ge}$ use
the original raw middle input
\[
 V(h)=\frac{w(u_+;h)}{\ell_0h_{\rm ax}(u_+)\rho(u_+)}
     =\frac{y(u_+;h)}{\varepsilon_+}>0,
 \quad h_{\rm ax}=-\delta p_d,\quad
 \varepsilon_+=\frac{\ell_0h_{\rm ax}(u_-)}{a^2J_-}.
 \tag{PCov2}
\]
All factors in this definition are actual and positive for a physical
coefficient fibre. Define the following subsets of that same carrier:
\[
\begin{aligned}
 G_\theta&=I_\theta^-\cap I_{\ell,\theta}
          \cap\{r_\ell\ge r_{\rm ge},\ \sigma\le\sigma_0\},\\
 D_\theta&=I_\theta^-\cap I_{\ell,\theta}
          \cap\{\sigma\ge\sigma_1,\ 0<V\le V_1\},\\
 E_\theta&=I_\theta^-\cap I_{\ell,\theta}
          \cap\{\sigma\ge\sigma_1,\ V\ge1\}.
\end{aligned}                                                    \tag{PCov3}
\]
The $V$ coordinate in the last two lines is well defined because
$\sigma_1/s_0>r_{\rm ge}$. These are connected, relatively closed
intervals, possibly empty, and
\[
 I_\theta^-=B_\theta\cup G_\theta\cup D_\theta\cup E_\theta.
 \tag{PCov4}
\]
Each of $G_\theta,D_\theta,E_\theta$ has the exact raw input buffer
\[
 0<q(h):=\frac{h-x_b}{s_{\rm in}}<q_{\rm ref}
       <\widehat p_*<\widehat p_0,
 \qquad s_{\rm in}=\Omega a^{2-\lambda_\gamma}U_{\rm in}>0.
 \tag{PCov5}
\]
The compact annular/growing overlap and the original product and
$V$ overlaps are on these actual inputs, with their equality dividers
included. Intersecting with any original closed raw-window or height
truncation preserves the cover; no divider is required to occur.
\item $G_\theta$ lies in the growing product domain, and
$D_\theta,E_\theta$ lie in the enlarged domains of the finite-$V$ and
fixed-exit maps constructed in
Proposition~\ref{p2:prop-finite-v-fe-actual-budget}.
The high-height interval is $I_\theta^+=I_\theta\cap\{h\ge h^0\}$.
Together with (PCov4) it covers the whole $I_\theta$, including every
non-root through input. On all these intervals the physical equation is
\[
 F(h-a_\theta)=h-a_\theta-f_\theta(m\mathfrak T_\theta(h)),
 \qquad \mathfrak T_\theta=d_\theta\circ T_\theta^\infty.
 \tag{PCov6}
\]
\end{enumerate}
The statements about actual raw maps concern $j>0$; the reciprocal
face $j=0$ has only the uniform one-sided normalized interpretation.
\end{proposition}
\begin{proof}
The independently constructed stopped and $h_L=0$ graphs bound the
through interval. Every intermediate graph reaches $u_c\in(-1,u_-)$ with
positive denominator. Scalar first variation makes $W_\theta$ strictly
increasing, with stopped limit zero.
For the upper comparison graph set $r_0=s_0L$, $r=-u$ and $\mathscr W=w/r^2$.
At its initial cut
\[
 \mathscr W(r_0)=\tfrac12+O_L(\tau b_\tau),\qquad
 b_\tau=-\log(2\tau).
\]
On $r_0\le r\le |u_c|$, direct substitution of the original polynomial
in the graph equation gives, on $1/4\le\mathscr W\le3/4$,
\[
 r\mathscr W_r=1-2\mathscr W+\mathscr E,\qquad
 |\mathscr E|\le C\{\kappa\tau+\kappa\tau/r\}
                 \le C\{\kappa\tau+\tau/L\}.             \tag{PCov7}
\]
Indeed $|\delta p_d(-r)|\le C\kappa\tau(r+a)$ and
$a/r\le x_{\max}/L$, while $K_d=1+O(\kappa\tau)$ and
$-L_d(-r)/r=1+O(\kappa\tau)$; the $E_dw$ term is smaller on this
corridor. The regular reciprocal factorization gives the same bounds
uniformly as $j\downarrow0$. Both boundary vectors point inward.
The logarithmic-radius equation contracts, so variation of constants
gives $\mathscr W(|u_c|)\to1/2$. Hence
$W_\theta(h^0)\to u_c^2/2$. The uniform endpoint margins put
$\overline{\mathcal J_c^{\sharp\sharp}}$ strictly inside this shooting
image after shrink. Strict first variation gives a unique through input
for each $W$; backward uniqueness identifies its shot and the two
already constructed raw continuations. This proves (i) independently
of the fixed-exit phase.

The same incoming first variation preserves the ordering of the
actual $h_L$ labels. Thus the central source interval is $[h_S,h^0]$.
The exact gate inverse and centre matching give the annular interval
$[h_B,h_A]$. The fixed orbit $h_L=-\delta_1$ has trace $h_C$ strictly
inside both intervals, because $-\delta_2<-\delta_1<0$ and its actual
gate labels lie strictly inside the annular bands. Consequently
\[
 [h_S,h^0]\cup[h_B,h_A]=[h_T,h^0].                       \tag{PCov8}
\]
The first fixed comparison alone supplies this overlap. The compact overlap also supplies an actual reference with $r_\ell=r_{\rm ge}$ and $r_b<2R_{\rm BD}$;
its trace $h_{\rm ge}$ satisfies $h_B<h_{\rm ge}$.
Every input failing (PCov8) is below $h_T\le h_B$, so the supplied full
original first-hit order gives $r_\ell>r_{\rm ge}$. The unified entrance
then gives (PCov5) on all inputs with $r_\ell\ge r_{\rm ge}$, including
first hits beyond the product range. Their derivative estimates are applied only within the growing product
domain.

For this actual local input $q$, the original positive cut is reached
on the stopped outer chord's local invariant. Its height is positive,
finite and strictly increasing in $h$. Formula (PCov2) is exactly the
original normalization $w=\ell_0h_{\rm ax}\rho V$, so
$V_h>0$, while the full first-hit order gives $\sigma_h<0$.
All sets in (PCov3) are therefore intervals and include their finite
level boundaries. If an input outside $B_\theta$ has
$\sigma\le\sigma_0$, it is in $G_\theta$. Otherwise
$\sigma>\sigma_0>\sigma_1$ and, because $V_1>1$, its positive $V$
satisfies $V\le V_1$ or $V\ge1$. It is in $D_\theta$ or
$E_\theta$. This proves (PCov4) for every through input, whether or not $F$ vanishes.

The compact band $r_\ell\in[r_{\rm ge},r_{\rm ge}^{\sharp}]$
is strictly inside the annular interval and, after the prescribed
$s_0$ bound, the growing product domain. The inequalities
$\sigma_1<\sigma_2<\sigma_0\le\bar\sigma<|u_c|/2$ retain the actual
product overlap. At $\sigma=\sigma_0$ the other passage region has a strict
$\sigma_1$ buffer, and at $V=1,V_1$ the last two regions overlap
because $1<V_1$. If a level is absent, the corresponding truncated interval is empty. The bounded endpoints and annular/growing overlap are included in these
intervals.

On $G_\theta$, $r_\ell\ge r_{\rm ge}>r_{\rm ge}^{\rm out}$ and
$\sigma\le\sigma_0\le\bar\sigma$, so the independently constructed
negative lip and positive chord identify the actual growing graph.
On $D_\theta\cup E_\theta$ the same raw buffer (PCov5) supplies the
original positive local cuts. The construction in Proposition~\ref{p2:prop-finite-v-fe-actual-budget}, using (i) as its
independent shooting geometry, identifies the bounded-$V$ early-cut
passage and the fixed-exit shooting map with precisely these orbits. In particular its action bound places the whole $E_\theta$
inside the enlarged inverse image, with $\Theta=0$ buffered. This last step is the only use of that proposition; the shooting
geometry in part (i) was constructed independently.

Finally the $h^0$ comparison graph has original source height
$Y_0\asymp b_\tau\to\infty$ and physical height
$H_\theta(h^0)=s_s^2Y_0\to0$. The original source dictionary gives
$|\alpha_0|+\alpha_1^2=s_s^2(1+p_s^2)$, with $p_s^2=O(\tau)$.
It follows that
\[
 K(|\alpha_0|+\alpha_1^2)<H_\theta(h^0)<z_{\rm cap}.
\]
Scalar height order puts every retained input at or above $h^0$ in
that same original high-height interval. The real enlarged $h_L=0$
source chord supplies its open overlap with the low carrier. Thus
$I_\theta^+$ includes the equality $h^0$ and all its non-root inputs;
the upper cap is common and fixed, rather than selected at a zero.
Original-time uniqueness gives (PCov6), with no raw gain inverted in
the domain construction.
\end{proof}

\section{Bounded source}

\begin{lemma}[Physical derivative estimates on the bounded-source intervals]
\label{p2:lem-bounded-source-budget}
Keep the original field, $L,h_L$, raw ports and enlarged windows, regular
maps and ordinary $C^3$ divided closing. Keep the unified tuple's literal
\[
 \begin{gathered}
 \delta_1=\min\{\delta_\sharp(L)/8,x_{\min}^2/8192,1/8192\},\\
 \delta_2=2\delta_1,\quad p_{\rm cmp}=\sqrt{2\delta_1},\quad r_a=p_{\rm cmp}/8,\quad p_a=r_a/16,\\
 p_-=p_a/2,\quad c_\pm=p_a/8.
 \end{gathered}
\]
The constant $\delta_\sharp(L)$ was fixed once, sufficiently small for
the bounded-source domain and relative estimate, before these gate
constants were defined. On the actual source word put
\[
 \beta=h_L/\tau-\log(2\tau),\quad b_\tau=-\log(2\tau),\quad
 \beta_{\min}=-\delta_2/\tau-\log(2\tau).
\]
There are fixed thresholds $B_-,B_+$ and constants $C_B,C_H$, uniform in
the original compact and reciprocal coefficient families, such that:
\begin{enumerate}
\item On every actual input with $\beta_{\min}\le\beta\le-B_-$,
the complete divided transition $\Pi$ has $\Pi''<0$ in the original
affine raw input, hence $F''=-m\Pi''>0$ for $m>0$.
\item On $-B_-\le\beta\le B_+$, one has $0<\Pi'\le C_B$.
\item On $B_+\le\beta\le b_\tau$, including $h_L=0$, and on a real
enlarged source chord through that orbit, one has $0<\Pi'\le C_H$.
The same bound holds on its admitted height-dominant continuation.
\end{enumerate}
All assertions hold on enlarged overlap intervals and at every nonclosed
through input. One final choice $m\max\{C_B,C_H\}<1/2$ gives $F'>1/2$
on the last two pieces. The input is always
$x=h-a_\theta$, and
\[
 \Pi(x)=\mathcal C_{\theta,m}(T_\theta(a_\theta+x)),\qquad
 F(x)=x-m\Pi(x)=x-f_\theta(m d_\theta(T_\theta(a_\theta+x))).
 \tag{BS1}
\]
\end{lemma}

\begin{proof}
Apply Proposition~\ref{bd:prop-bounded-action} to the same sections
and divided closing. Its estimates
\eqref{bd:relative-negative-curvature}--\eqref{bd:relative-raw-curvature}
give the first assertion, and \eqref{bd:compact-raw-variation}
gives the second.
The independent frozen windows, $L_f=8$, $R_f=5$ and polynomial
margin $3/4$ are supplied by Lemma~\ref{bd:lem-source-ranges};
both exact raw gains and their common-power cancellation are retained
in (BS1). It remains to check the literal gates and the upper chord.

The final $\delta_\sharp(L)$ is chosen there for both the domain and
relative curvature estimate, after $B_-$ and before the displayed literal
$\delta_1,\delta_2$. Since $\delta_2\le\delta_\sharp/4$,
the central block has enlarged lower chords down to
$h_L=-3\delta_1>-\delta_\sharp$. Compact labels
$[-B_--1,B_++1]$ and positive labels beginning below $B_+$ give actual
overlaps by monotone incoming first variation. A final $\tau$ shrink
makes $\beta_{\min}<-B_-$ and $b_\tau>B_++2$.

For positive labels choose the signed high-height constants
$K,z_*,C_H$ first, then $B_+$ large. Equation~\eqref{bd:positive-height}
gives, through $\beta=b_\tau$ and the enlarged chord below,
\[
 Y_0>K(1+p_s^2),\qquad
 H=s_s^2Y_0\le C_L\tau^3\sigma^4j^2(1+b_\tau)<z_*.
 \tag{BS10}
\]
To realize the upper chord, freeze $B=b_\tau$ in the original
positive-energy construction. For $\beta/B\in[3/4,3/2]$, its incoming
energy $\beta/B+B^{-1}\log(2/B)$ stays strictly inside $[1/4,3]$.
The energy defect is $O_L(B^{-1})$ and its variation integral tends
uniformly to zero; the same first-contact proof and positive inverse
therefore construct this entire chord. It has
$Y_0=\beta-\log\beta+O_L(1)$ and
$\eta(\pm L)=L^2/2+O_L(\tau b_\tau)$ at the unchanged cuts.
The original quadratic corridor and raw-window margins continue both
sides. In particular $0\le h_L\le\tau b_\tau/2$ is a real enlargement
above the $h_L=0$ orbit.

By (BS10), the signed high-height bound used in
Proposition~\ref{bd:prop-bounded-action}, namely
Corollary~\ref{rew:cor-high-derivative-signed}, applies throughout
this chord and its continuation
$K(|\alpha_0|+\alpha_1^2)\le H<z_*$, including the equality face.
It gives the third assertion, with the height scale frozen at each base point.

The full transported output, including reference values, has a common
bound $M_{\rm out}$ on the original section windows. Choose $m$ last
with $m\le m_0$, $mM_{\rm out}<a_{\rm cl}/2$ and
$m\max(C_B,C_H)<1/2$. Then (BS1) gives $F'>1/2$ on the compact,
positive and high pieces, while the negative curvature gives $F''>0$
for $m>0$. The enlarged estimates include threshold equalities and
every nonclosed input. At $m=0$, $F(x)=x$; collapsed sections denote
only normalized one-sided limits.
\end{proof}

\section{Growing exit}

\paragraph{The growing displacement and its curvature.}
The action derivative of the matched map is known, but the finite
closing is evaluated at its full translated output. We first bound that
output and its first two action derivatives, then use the logarithmic
slope to prove convexity in the original affine input.
All phase derivatives fix every coefficient. Use the already constructed
joint auxiliary tuple, the actual first lip $r=r_\ell$, and its original
enlarged domain
\[
 \widehat r_{\rm ge}^{\rm out}\le r\le\bar\sigma/s_0,\qquad
 s_0=\kappa\sqrt B,\qquad 0<\sigma_0<\bar\sigma<1/2 .
 \tag{GC1}
\]
The counted interval has $r\ge\widehat r_{\rm ge}$ and $s_0r\le\sigma_0$.
Write the actual affine raw input as
\[
 h=x_b+s_{\rm in}P,\qquad
 s_{\rm in}=\Omega a^{2-\lambda_\gamma}U_{\rm in}>0,\qquad
 \Omega=r_s^{2-\lambda_\gamma},\qquad U_{\rm in}^{\pm1}\le C .
 \tag{GC2}
\]
The joint entrance construction supplies $0<P<\widehat p_*/2$.
Set
\[
 X=\Phi_\theta(P)/\Phi_\theta(q_{\rm ref}),\qquad
 \mathcal E=X\partial_X=\mathfrak j(P)\partial_P,\qquad
 \mathfrak j=\Phi_\theta/\Phi_\theta'>0,\qquad
 e_s=\tau/\chi ,
 \tag{GC3}
\]
with the denominator held fixed. Let
\[
 \mathcal P(r)=P_s(-r)/B_\tau(-r),\qquad
 \mathfrak d=\mathcal P/r,\qquad \mathcal D=\mathfrak d\partial_r .
\]
The polynomial estimates and the matched action derivative give
\[
 \begin{gathered}
 \mathcal P\asymp1+r,\quad \mathcal P'>0,\quad
 \mathfrak d^{\pm1}\le C,\quad |\mathcal D\log\mathfrak d|\le C,
 \quad e_s^{\pm1}\le C,\quad \mathfrak j'>3/4,\\
 -\mathcal Er=e_s\mathfrak d\,n,\qquad n=1+\zeta,\qquad
 n\ge1/2,\qquad |\zeta|+|\mathcal E\zeta|\le C\tau .
 \end{gathered}
 \tag{GC4}
\]
Here $\mathcal E$ acts on the actual graph, including the negative
variable in the complete quotient row.

Retain the full actual output
\[
 \begin{gathered}
 T(r)=\Xi(W_\theta(r))
      =t_\dagger-a_{\rm grow}\mathcal B(r),\qquad
 a_{\rm grow}=\Omega s_0^2U_{\rm out}>0,\qquad
 \mathcal B'=rG,\quad G>0,\\
 W_\theta'=-s_0^2r\Gamma_\theta,\qquad
 G=\Gamma_\theta\,\Xi'(W_\theta)/(\Omega U_{\rm out}) .
 \end{gathered}
 \tag{GC5}
\]
Thus $T$ includes its coefficient-dependent translation, original
moving endpoint, tails and regular output section map.
The actual output lies in a fixed compact physical section window $J_T$.
The first action derivative of the landing, composed with this same
$\Xi$, satisfies
\[
 \begin{split}
 |\mathcal D\log G|
 &\le C\left\{
 \frac{\tau}{\sqrt B\,r^2}
       \left(1+\left|\log\frac{\tau}{\sqrt B\,r}\right|\right)
       +s_0\kappa\tau+s_0^2\mathcal P(r)\right\}\\
 &\le C\{\tau(1+|\log\tau|)+s_0\}.
 \end{split}
 \tag{GC6}
\]
The ordinary finite closing has $f_\theta(0)=0$, $f_\theta'>c_f>0$,
and uniform ordinary $C^3$ bounds on a buffered interval containing
$mJ_T$ for $0\le m\le\bar m$. It is the original closing after the
output map $d_\theta$ already retained in $T$.
The physical affine displacement is exactly
\[
 F_\theta(h)=h-a_\theta-f_\theta(mT_\theta(h)),\qquad
 T_\theta(h)=T(r(P(h))).
 \tag{GC7}
\]

\begin{proposition}[Convexity and zero count in the growing range]
\label{p2:prop-growing-physical-count}
On the whole counted growing interval, there are common constants
$C_T,C_{\rm cl}$ and a multiplier bound $m_0>0$, selected after the
independent coefficient bounds, such that
\[
 |T|\le C_T,\qquad 0<\mathcal ET\le C_T,\qquad
 |\mathcal E^2T|\le C_T .
 \tag{GC8}
\]
For every actual through input in that interval, whether or not
$F_\theta=0$,
\[
 \begin{gathered}
 \left|\mathcal E^j\{f_\theta(mT)-f_\theta(0)\}\right|
       \le C_{\rm cl}m,\\
 \left|\mathcal E^j\{\log f_\theta'(mT)-\log f_\theta'(0)\}\right|
       \le C_{\rm cl}m,\qquad j=0,1,2,\quad 0\le m\le m_0 .
 \end{gathered}
 \tag{GC9}
\]
For $0<m\le m_0$, the complete physical displacement is strictly convex:
\[
 F_\theta''(h)>0\qquad\hbox{at every point of the growing interval}.
 \tag{GC10}
\]
Consequently it has at most one critical point, that point is simple
and upward, and it has at most two isolated zeros counted with
multiplicity. At $m=0$, $F_\theta'(h)=1$.
The constants are uniform on the full prescribed passive coefficient
family and both coefficient graphs, with the reciprocal estimates
uniform for physical $j>0$ approaching its one-sided boundary.
\end{proposition}

\begin{proof}
\emph{Bounds on the full output and its action derivatives.}
Put $L(r)=-T_r=a_{\rm grow}rG(r)>0$ and
\[
 d_r=\min\{1,(\widehat r_{\rm ge}
              -\widehat r_{\rm ge}^{\rm out})/2\}>0 .
\]
The original lower buffer places every $[r-d_r,r]$ in (GC1).
Since $\mathfrak d^{-1}$ is bounded, (GC6) gives
\[
 |\partial_r\log L|
 \le r^{-1}+\mathfrak d^{-1}|\mathcal D\log G|\le C_L .
\]
Let $D_T$ be the diameter of $J_T$. Using the actual output values,
\[
 D_T\ge T(r-d_r)-T(r)
      =\int_{r-d_r}^rL(v)\,dv
      \ge d_re^{-C_Ld_r}L(r).
 \tag{GC11}
\]
Thus $L\le C$ without dividing by $a_{\rm grow}$, $\Omega$ or a raw gain.
Equation (GC4) now gives
\[
 \mathcal ET=e_s\mathfrak d\,Ln>0,\qquad
 \mathcal E\log(\mathcal ET)
 =-e_sn\{\mathcal D\log\mathfrak d+\mathcal D\log L\}
       +\frac{\mathcal E\zeta}{n},\qquad
 \mathcal D\log L=\mathfrak d/r+\mathcal D\log G .
 \tag{GC12}
\]
Every term on the right is bounded. Hence $\mathcal ET\le C$ and
$|\mathcal E^2T|\le C$; $|T|$ is bounded by the fixed compact $J_T$.
The reference value of the output is included in this bound, so (GC8)
applies at the actual argument of the closing map.

\emph{Composition with the finite closing.}
The full translated $T$ and its first two action rows satisfy (GC8)
on the buffered actual image. The ordinary $C^3$ closing has positive
derivative units there. Lemma~\ref{phys:lem-c3-adapters}, specifically
\eqref{phys:eq-absolute-closing-two}, therefore gives (GC9) for all
through inputs, with constants fixed before $m$.

\emph{Strict convexity in the actual affine input.}
Before finite closing, the exact logarithmic slope is
\[
 \begin{split}
 H_0
 &:=\mathcal E\log
       \frac{\mathfrak j(P)}{e_sG(r)\mathcal P(r)n(P,r)}\\
 &=\mathfrak j'(P)+e_sn
       \left\{\frac{\mathcal P'(r)}r+\mathcal D\log G(r)\right\}
       -\frac{\mathcal E\zeta}{n}.
 \end{split}
 \tag{GC14}
\]
This identity follows from $\mathcal Er=-e_s\mathfrak d n$;
the derivative of the matched quotient is included.
By (GC4), (GC6) and $\mathcal P'>0$,
\[
 H_0\ge 3/4-C\{\tau(1+|\log\tau|)+s_0\}>1/2
 \tag{GC15}
\]
after common independent $\tau,\kappa$ bounds.
Here $s_0=\kappa\sqrt B$, with $B$ a positive unit.
The original passive coefficient set and joint auxiliary tuple are
unchanged.

Since $\mathcal ET=\mathfrak j T_P$ and (GC5) gives
$T_P=a_{\rm grow}e_sG\mathcal P n/\mathfrak j>0$, define
\[
 \mathscr R=\frac{s_{\rm in}}{mf_\theta'(mT)T_P}>0 .
\]
Every factor in its numerator, $a_{\rm grow}$ and $e_s$ is
coefficient-only. Its complete logarithmic row is therefore
\[
 H_m:=\mathcal E\log\mathscr R
    =H_0-m\frac{f_\theta''(mT)}{f_\theta'(mT)}\mathcal ET .
 \tag{GC16}
\]
The absolute output estimate (GC8) makes the closing term at most
$C_m m$, with a common $C_m$ fixed before the multiplier.
Choose last
$m_0\le\min\{\bar m,1,1/[4(1+C_m)]\}$. Then $H_m>1/4$.

Restore the physical input using (GC2). Direct differentiation of the
same displacement (GC7) gives at every through point
\[
 F_\theta'(h)=1-\frac1{\mathscr R},\qquad
 F_\theta''(h)=\frac{\mathscr R_h}{\mathscr R^2}
       =\frac{H_m}{s_{\rm in}\mathfrak j(P)\mathscr R}>0 .
 \tag{GC17}
\]
The two raw gains remain in this expression, so its sign holds even
when either gain is arbitrarily small.
Translation to $h-a_\theta$ leaves these derivatives unchanged.

Lemma~\ref{sharp:upper-critical-count} now gives the critical and
zero counts, with multiplicity; (GC17) excludes an identity interval.

The bound (GC6) uses $r\ge R_0$ and
$s_0^2\mathcal P(r)\le Cs_0$; its logarithm contains no $j$.
The bounds on $\kappa,\tau$ are independent, uniformly on the
prescribed compact $j$ ranges.
\end{proof}

\section{Finite height and fixed exit}
\subsection{The last two passage regions and their physical curvature}

Bounded amplified height is described by two small-power section maps.
For larger height we instead use the action coordinate on a fixed
negative section. We construct both maps on overlapping domains before
comparing their derivatives in the physical input.

Retain the single joint auxiliary choice, its original physical sections
and nested windows, the weak cuts, and the original coefficient family.
Let \(I_\ell\) be the positive-height first-lip domain in the original
through interval. Write
\[
 P={\frac{h-x_b}{s_{\rm in}}},\qquad
 s_{\rm in}=\Omega a^{2-\lambda_\gamma}U_{\rm in}>0,\qquad
 \Omega=r_s^{\,2-\lambda_\gamma},
 \qquad q_s={\frac{\beta}{\chi}}=\kappa Q_d ,
 \tag{FVFE1}
\]
where \(h\) is actual incoming height, \(Q_d,U_{\rm in}\) are positive
units, and no positive lower bound for either raw first gain is imposed.
The variable \(P\) is the canonical raw input of the joint entrance
construction; it is distinct from the saddle ratio \(q_s\).
Use the unchanged raw middle coordinate
\[
 V={\frac{y_+(h)}{\varepsilon_+}},\qquad
 \varepsilon_+={\frac{\ell_0h_-}{a^2J_m(u_-)}},\qquad
 \sigma=s_0r_\ell .
 \tag{FVFE2}
\]
Here \(y_+\) is the actual height \(w/a^2\) at \(u=u_+\).
Every phase derivative fixes all coefficients and all gains.

\begin{proposition}[Physical passage and convexity in the last two ranges]
\label{p2:prop-finite-v-fe-actual-budget}
For the joint auxiliary choice, followed by independent active
coefficient bounds, put
\[
 J_{34}=I_\ell\cap\{\sigma\ge\sigma_1\},\qquad
 D=J_{34}\cap\{0<V\le V_1\},\qquad
 E=J_{34}\cap\{V\ge1\},\qquad V_1>1 .
 \tag{FVFE3}
\]
These are intervals in the same actual incoming height, with every
displayed divider equality included. Each lies in the prescribed
enlarged domain of its corresponding actual transition, and
\(D\cup E=J_{34}\). Every input in these sets, whether
closed or nonclosed, has its original raw trace in
\(0<P<q_{\rm ref}<\widehat p_*<\widehat p_0\).
The finite-\(V\) map uses the original early negative cut on
the fixed enlarged interval \(0<V\le V_1+1\). For larger amplified
heights, the independently constructed fixed-exit map uses the buffered
shooting chord on \(u=u_c\). The two
constructions agree on \(1\le V\le V_1\) by original-time
uniqueness, so their union is the same physical transition on \(J_{34}\).

For the original return write
\[
 \mathcal F(h):=F(h-a_\theta)
       =h-a_\theta-f_\theta(mT(h)).
 \tag{FVFE4}
\]
The map \(T\) includes the actual negative output adapter. For every
\(0<m\le m_0\),
\[
 \mathcal F''(h)>0\qquad(h\in J_{34}).
 \tag{FVFE5}
\]
Thus the whole combined interval has at most one simple upward
critical point and at most two isolated zeros with multiplicity.
At $m=0$, $\mathcal F'=1$. The whole carrier is counted separately
in Proposition~\ref{p2:prop-physical-degree-count}.
\end{proposition}

\begin{proof}
\emph{The shooting chord and its phase coordinate.}
Part (i) of Proposition~\ref{p2:prop-unified-carrier-cover} realizes
$\overline{\mathcal J_c^{\sharp\sharp}}\Subset(0,u_c^2/2)$ as
unique through orbits, identified with the backward shots by uniqueness.
Proposition~\ref{fx:prop-inverse-phase} applies to this entire buffered
chord: its scalar first-exit construction is independent of the
positive raw-section dictionary. Its extracted phase
$\Psi_\theta(W)$ converges in $C^2$ to
$\Gamma_c\circ\mathcal W^{-1}$, whose derivative is
$1/H_{0,c}(r)>0$. Hence its increasing inverse
$W=\mathcal H_\theta(\Theta)$ has bounded first two derivatives
and a positive derivative unit on
\[
 \Theta_-\le\Theta\le\Gamma_- -\Delta,\qquad \Theta_-<0,
 \tag{FVFE8}
\]
strictly inside its inverse image. The negative buffer comes from the
original enlarged chord around $u_-$; the reciprocal functions give
the same uniform gaps.

\emph{Continuation of the shooting chord to the raw input.}
The exact positive invariant, including its local unit functions,
is
\[
 \Phi_\theta(P)=\{z_+y_+(P)\}^{\chi},\qquad
 \Phi_\theta(P)=P U_\theta(P),\qquad
 \Phi_\theta'>c,\qquad U_\theta^{\pm1}\le C.
 \tag{FVFE9}
\]
Use the reference $P_{\rm ref}=a_{\rm in}$ of
\eqref{deg:degree-fixed-exit-gap}:
\[
 \Phi_\theta(P_{\rm ref})=(z_+\beta\varepsilon_+)^\chi,\qquad
 q_s\log\Phi_\theta(P_{\rm ref})=-\Gamma_-+o(1).
 \tag{FVFE10}
\]
The exact identity \eqref{p2:reference-telescope} proves this limit
on both graphs: on the reciprocal graph use its prescribed $d,j$
in place of $\mathcal Rc,1/c$. All factors remain regular at $j=0$,
so its error bound is unchanged and no $\beta\log j$ occurs.
This reference corresponds to $V=\beta$. On the compact shooting chord the limiting
gap from \(\Gamma_-\) has a fixed positive minimum \(d_{\rm fx}\).
The scalar shooting estimate and (FVFE10) give, after shrink,
\[
 \beta\log V\le\Gamma_- -d_{\rm fx}/2,\qquad
 \beta\log y_+\le-d_{\rm fx}/3 .
 \tag{FVFE11}
\]
Since \(\beta\le C\chi\), this implies \(y_+<\chi^2\).
Thus (FVFE9) can be inverted on the original invariant chord, and
\[
 \Phi_\theta(P)\le C e^{-d_{\rm fx}/(3q_s)},\qquad
 0<P\le e^{-d_{\rm fx}/(4q_s)}<\widehat p_0
 \tag{FVFE12}
\]
after a further independent \(\kappa\) bound.
In particular the physical distance from the stopped trace satisfies
\(\sup(h-x_b)\le C e^{-d_{\rm fx}/(4q_s)}\to0\) on the
entire shooting chord, since \(s_{\rm in}\le C\).
The two local points have the same invariant in the same larger
coordinate box. Positivity and monotonicity along that invariant level
give the intervening orbit segment; the independently actual raw
continuation identifies its raw trace with the shot through orbit.
This places the full chord within the independently selected fixed cutoff
$\widehat p_0$, before any zero selection.

\emph{The bounded-\(V\) passage.}
Lemma~\ref{lem:actual-dulac-matching} independently continues the
entire larger interval $0\le V\le V_1+1$ to the early cut
$u_0=-1+\eta_0>u_-$. Equations~\eqref{eq:early-saddle-smallness}
and~\eqref{eq:early-middle-identity} give its positive invariant
at most $Ce^{-c_0/(2q_s)}$. The inversion used in (FVFE12) places the whole interval strictly
inside the raw cutoff.
The assigned set \(D\) is its restriction to \(0<V\le V_1\).
On the negative early cut, the lemma's exact invariant is
\[
 I_-(w)=\alpha_\theta R,\qquad
 R=\beta^{-\beta}C(v)^\beta,\qquad
 v=\widehat V=V+o_{C^3}(\chi).                 \tag{FVFE14}
\]
Its fixed cut choice puts the full $R\le1+o(1)$ interval strictly
inside the outer positive-height chord. At $R=0$ the output is the
stopped height $w_->0$. Thus both actual tails, ordinary $C^3$
observers and finite closing are on their original buffered domains,
with their exact values and positive first gains retained.

\emph{The domains of the two passage maps.}
The joint entrance supplies the actual product reference
\(r_\ell=\sigma_1/s_0\) after the active shrink.
Direct landing followed by the independent inverse shooting phase
gives
\[
 \Theta(h_{\sigma_1})\longrightarrow\Gamma_c(-\sigma_1),
 \qquad
 \Gamma_- -\Gamma_c(-\sigma)
      =\log{\frac{c+\sigma}{c(1-\sigma)}} .
 \tag{FVFE15}
\]
On the reciprocal chart the gap is
\(\log((1+j\sigma)/(1-\sigma))\), uniformly positive on the fixed
product buffer. The retained choice of \(\Delta\) and its strict
buffer imply
\[
 0<g_0\le\Theta(h_{\sigma_1})\le\Gamma_- -2\Delta .
 \tag{FVFE16}
\]
The actual orders \(r_{\ell,h}<0\) and \(V_h>0\) imply that every
member of \(E\) has \(0\le\Theta=\beta\log V\le\Theta(h_{\sigma_1})\).
Thus the entire fourth range lies in the domain (FVFE8) of the inverse
phase map, including inputs whose orbits are not closed. The equality \(V=1\) has the
negative-action buffer there. Since
\(\mathcal H_\theta([0,\Gamma_- -2\Delta])
\Subset\mathcal J_c^\sharp\Subset\mathcal J_c^{\sharp\sharp}\)
uniformly after shrink, its whole closed output-height subset has
a strict actual shooting buffer.
Decrease the independent coefficient bounds further so that
\(\beta\log(V_1+1)<g_0/2\). Then
\(V(h_{\sigma_1})=\exp(\Theta(h_{\sigma_1})/\beta)>V_1+1\),
so the third range retains the full buffered finite-\(V\) passage.
The two constructions identify the same raw inputs and outputs on \(1\le V\le V_1\) by original-time uniqueness.
No physical input is relabelled at \(V=1\).
Their actual \(\sigma\ge\sigma_1\) domain has
\(r_\ell\ge\sigma_1/s_0>r_{\rm ge}\), so the joint entrance's full
first-hit order supplies the sharper raw bound
\(P<q_{\rm ref}<\widehat p_*<\widehat p_0\).

\emph{Physical convexity for bounded \(V\).}
For \(m>0\), the actual-section construction gives exactly
\[
 \begin{gathered}
 h=x_b+P_*K(v),\qquad
 f_\theta(mT(h))=b_*+a_*H(v),\\
 K(v)=\Psi(v^\chi),\quad
 H(v)=D(\beta^{-\beta}C(v)^\beta),\qquad P_*,a_*>0 .
 \end{gathered}\tag{FVFE17}
\]
Here all regular tails and finite closing are already included.
Put \(e_C=vC'/C>0\). Lemma~\ref{lem:actual-dulac-matching}
supplies the stronger normalized estimates below, with \(\Sigma_+\)
the coefficient-only positive invariant scale in its construction:
\[
 \begin{gathered}
 e_C=1+O(\epsilon_\theta),\qquad
 \left|{\frac{ve_C'}{e_C}}\right|\le C\epsilon_\theta,\qquad
 |D''/D'|\le C,\\
 \sup_t|t\Psi''(t)/\Psi'(t)|\le C\Sigma_+,\qquad
 \epsilon_\theta\le C\beta^{-N}e^{-c_0/(3\beta)},\qquad
 \Sigma_+\le Ce^{-c_0/(2q_s)} .
 \end{gathered}
\]
The exact Wronskian row, with \(R=\beta^{-\beta}C(v)^\beta\), is
\[
 \begin{split}
 \mathcal Q(v)
 &=v\left({\frac{K''}{K'}}-{\frac{H''}{H'}}\right)\\
 &=\chi-\beta e_C-{\frac{ve_C'}{e_C}}
       +\chi v^\chi{\frac{\Psi''}{\Psi'}}
       -\beta e_C R{\frac{D''}{D'}}
 \ge\chi/2>0 .
 \end{split}
 \tag{FVFE18}
\]
On this interval \(R\le1+o(1)\). With
\(q_s=\beta/\chi=\kappa Q_d\longrightarrow0\),
the displayed estimates give the relative bound
\[
 \frac{\mathcal Q(v)}{\chi}
 \ge1-Cq_s-C\frac{\epsilon_\theta}{\chi}-C\Sigma_+
 \ge\frac12 .
\]
Indeed, \(\beta=q_s\chi\) implies
\[
 \frac{\epsilon_\theta}{\chi}
 \le Cq_s\beta^{-N-1}e^{-c_0/(3\beta)}\longrightarrow0.
\]
The input-section error has weight \(\chi\Sigma_+\), whereas
the bounded output and closing error has weight \(\beta\).
Thus independent bounds on \(\kappa,\tau\) give this sign for
every \(v>0\), without extending the small-power compositions
in ordinary \(C^3\) to the excluded stopped endpoint.
Since \(dh/dv=P_*K'>0\), differentiation of the exact physical
displacement, without multiplying it by a nonconstant factor, yields
\[
 \mathcal F''(h)=
 {\frac{a_*}{P_*^2}}{\frac{H'}{(K')^2}}{\frac{\mathcal Q(v)}{v}}>0 .
 \tag{FVFE19}
\]

\emph{Physical convexity on the fixed-exit interval.}
The assigned set \(E\) has
\(0\le\Theta\le\Theta(h_{\sigma_1})\), strictly inside the
larger buffered action interval (FVFE8).
Use the common actual map \eqref{deg:fixed-exit-actual-output},
$\Xi(W)=t_b+a_{\rm out}D^n(W)$, $a_{\rm out}=\Omega U_{\rm out}>0$,
on its whole buffered height chord. Its ordinary $C^3$ normalized
bounds and $0<\Omega\le C$ retain the full output in the original
raw and closing windows. For the fixed closing buffer $a_{\rm cl}>0$,
impose on the whole chord
\[
 m_0\left\{|t_b|+
 a_{\rm out}\sup_{\mathcal J_c^{\sharp\sharp}}|D^n|\right\}
 <a_{\rm cl}/2 .
 \tag{FVFE20a}
\]
The braces are uniformly bounded on the original output window,
so this selects $m$ independently of the first gain.

Put \(\widehat H=D^n\circ\mathcal H_\theta\). The exact raw formulas are
\[
 T(h)=t_b+a_{\rm out}\widehat H(\Theta(P)),\qquad
 \Theta(P)=\beta\log\beta+
       q_s\log{\frac{\Phi_\theta(P)}{\Phi_\theta(P_{\rm ref})}}.
 \tag{FVFE21}
\]
Writing \(p=h-x_b=s_{\rm in}P\), retain the raw gain in the two rows
\[
 0<pT_h=
 a_{\rm out}\widehat H' q_s{\frac{P\Phi_\theta'}{\Phi_\theta}}
       \le C a_{\rm out}q_s ,
 \tag{FVFE22}
\]
\[
 p{\frac{T_{hh}}{T_h}}
 =P{\frac{\Phi_\theta''}{\Phi_\theta'}}
       -P{\frac{\Phi_\theta'}{\Phi_\theta}}
       +q_sP{\frac{\Phi_\theta'}{\Phi_\theta}}
                    {\frac{\widehat H''}{\widehat H'}}
 =-1+O(P+q_s)\le-\tfrac12 .
 \tag{FVFE23}
\]
The ordinary $C^3$ closing has derivative units and bounded
pre-Schwarzian on the whole sampled interval. (FVFE22)--(FVFE23) give
\[
 p\left\{{\frac{T_{hh}}{T_h}}
       +m{\frac{f_\theta''(mT)}{f_\theta'(mT)}}T_h\right\}
 \le-\tfrac12+Cma_{\rm out}q_s<-\tfrac14 .
 \tag{FVFE24}
\]
Therefore
\[
 \mathcal F''(h)=
 -mf_\theta'(mT)T_h
 \left\{{\frac{T_{hh}}{T_h}}
       +m{\frac{f_\theta''(mT)}{f_\theta'(mT)}}T_h\right\}>0 .
 \tag{FVFE25}
\]
No lower bound for \(s_{\rm in}\), \(\Omega\), or \(T_h\) was used.

The positive input order makes \(J_{34}\) one interval.
The two convexity estimates concern the identical physical \(F\)
on an overlapping cover, so (FVFE5) holds on their union.
Lemma~\ref{sharp:upper-critical-count} gives the critical and zero
counts with multiplicity, also for every finite collection approaching
the excluded stopped boundary.
The endpoints at the phase boundaries lie inside the enlarged domains
of their respective passage maps.
At \(m=0\) the original completed displacement has derivative one.
The combined reference logarithms, normalized section maps and
first-cocycle powers above are regular at the reciprocal face; none introduces a
\(\beta\log j\) or \(\tau\log j\) restriction. Actual physical maps
are asserted for \(j>0\), with \(j=0\) a one-sided normalized limit.
\end{proof}

\section{Counting zeros on the complete degree-loss interval}

Each local calculation either gives strict convexity of the physical
displacement or excludes critical points by a positive derivative.
Because the local maps agree on overlapping intervals, the same
alternative holds at every point of the connected through interval.

\begin{proposition}[Physical counting on the degree-loss carrier]
\label{p2:prop-physical-degree-count}
Fix the compact positive-unit ranges of the degree-loss graph, or
the reciprocal graph with $0<j\le j_0$. Use the joint auxiliary
choice of the growing entrance and its actual gates. There are
independent active coefficient bounds, a common raw-height cap and
a last finite-closing multiplier bound with the following property.
On the resulting connected interval $I_\theta$ of actual through
inputs on the original section, the complete affine displacement
\[
 F_\theta(x)=x-f_\theta\bigl(m\mathfrak T_\theta(a_\theta+x)\bigr),
 \qquad x=h-a_\theta,
 \tag{PC1}
\]
has at most one critical point. Every critical point is a simple
upward crossing of $F_\theta'$. Consequently it has at most two
isolated zeros, counted with multiplicity, and no identity interval.
The interval $I_\theta$ consists of the same through word in the
original nested windows, with positive passage denominator and
$0<z(\gamma_s)<z_{\rm cap}$ for one fixed physical height cap.
The assertion holds at every actual input of $I_\theta$, without
assuming its orbit is closed. The original field, marked sections,
nested raw windows, $L,h_L$ and closing map define this same return.
Here $\mathfrak T_\theta=d_\theta\circ T_\theta^\infty$ is the
original transported output and $f_\theta=C_\theta-C_\theta(0)$.
Thus $f_\theta(0)=0$ and at $m=0$ the displacement is $F_\theta(x)=x$.
No physical passage at $j=0$ or at a stopped boundary is asserted.
\end{proposition}

\begin{proof}
Proposition~\ref{p2:prop-unified-carrier-cover} supplies the closed
low cover $B_\theta\cup G_\theta\cup D_\theta\cup E_\theta$ and
its overlapping high continuation on the connected actual interval
$I_\theta$. It includes every divider equality and every nonroot input;
only the stopped endpoint is excluded.
All restrictions have the same affine displacement (PC1).

On the annular part of $B_\theta$,
Lemma~\ref{bd:lem-annulus} gives $\Pi''<0$ after the ordinary
$C^3$ divided closing is included. Its principal margin is
$\lambda_1/(2x_{\max}R_{\rm BD})$; the independent active bounds
make the errors smaller than half this margin.
Lemma~\ref{p2:lem-bounded-source-budget} gives the same sign on
negative source inputs and $F'>1/2$ on the compact, positive and
high source pieces. Propositions~\ref{p2:prop-growing-physical-count}
and~\ref{p2:prop-finite-v-fe-actual-budget} give $F''>0$ on
$G_\theta$ and on the whole combined $D_\theta\cup E_\theta$.
These signs are in the original affine input and include both exact
raw gains and the translated output inside the closing map.

Take the finite intersection of the independent active bounds, then
one multiplier bound for all derivative estimates and the full absolute
output envelope in the original closing buffer. Every critical point
of (PC1) consequently has $F''>0$, including those on the dividers.
Lemma~\ref{sharp:upper-critical-count} gives at most one simple
upward critical point and the two-zero bound with multiplicity.
An identity interval would contain a critical point with $F''=0$.
Section~\ref{sec:common-cover-conclusion} assembles the coefficient
cover and Hausdorff capture.
\end{proof}

\endgroup

The growing-action estimates can now be applied on the physical
interval. They control the matched quotient, which determines
curvature before finite closing:
\begin{equation}\label{deg:degree-growing-elasticity}
 \mathcal E\log\frac{\mathfrak j(P)}
    {e_sG(r)\mathcal P(r)n(P,r)}>\frac12,
 \qquad \mathcal E=X\partial_X.
\end{equation}
This is \textnormal{(GC14)--(GC15)}. To compose with the finite
closing, we also need the full output and its first two action derivatives:
\begin{equation}\label{deg:growing-actual-output-action}
 |T|+|\mathcal E T|+|\mathcal E^2T|\le C,
 \qquad \mathcal E=X\partial_X,\quad \mathcal E T>0.
\end{equation}
Estimate \textnormal{(GC8)} gives this bound on the buffered growing
image, and \textnormal{(GC9)} transfers it through the finite closing.
